% !TeX root = ../../Book4.tex
% 纲要标题：## 第13章*　Hochschild 同调与代数形变
% 纲要标签：b4:ch:12
% 纲要位置：计划纲要.md，第 3529 行。

\OptionalChapter{Hochschild 同调与代数形变}{b4:ch:12}
\BookChapterNumber{b4:ch:12}

\begin{chapterabstract}
	本章以包络代数和 bar 消解计算 Hochschild 同调与上同调，解释中心、导子和乘法形变。在交换代数上，星积的一阶交换子产生 Poisson 括号；在 Hochschild 链上加入 Connes 算子，则得到循环同调及其与 Hochschild 同调相连的长正合列。矩阵代数、多项式代数及平方零代数提供具体计算。
\end{chapterabstract}

\begin{learninggoals}
	\begin{itemize}
		\item 准确区分双模对应的左 \(A^e\) 作用与计算 Tor 所用的右作用。
		\item 从 bar 消解推导 Hochschild 微分，说明底层可缩与双模可缩的区别。
		\item 计算低阶同调和上同调，并写清形变分类中的单位条件及等价关系。
		\item 使用杯积和插入符号，计算特定形变的三次障碍类。
		\item 用矩阵单位证明矩阵代数的同调不变性，比较多项式的 Koszul 计算与光滑情形的局部证明。
		\item 区分星积的一阶结合条件与 Poisson 括号的 Jacobi 恒等式，并计算显式的形式星积。
		\item 核验 Connes 算子与 Hochschild 边界的相容关系，由循环复形的短正合列推导 SBI 长正合列并进行计算。
	\end{itemize}
\end{learninggoals}

令 \(t\) 为形式参数，给结合代数 \(A\) 的乘法加上一阶修正，写成 \(a*b=ab+t\cdot\phi(a,b)\pmod{t^2}\)。比较 \((a*b)*c\) 与 \(a*(b*c)\) 的 \(t\) 项，便得到
\[
	\phi(a,b)c+\phi(ab,c)=a\phi(b,c)+\phi(a,bc).
\]
因此 \(\phi\) 不能任意选取。若只改变元素的记法而不改变所得代数，这个修正项又会发生可控制的变化。我们要把「满足结合律」与「只是换了记法」分别写成余圈与余边条件，再看怎样用 Ext、Tor 和显式复形计算它们。

\ProseParagraphBreak
本章使用\cref{b4:thm:resolution-comparison}以及\cref{b4:thm:ext-balanced,b4:thm:tor-balanced}；多项式算例还使用\cref{b4:prop:koszul-resolution}，循环同调的比较使用\cref{b4:thm:cohomology-long-exact}的链指标形式。全章底域为 \(k\)，代数保持结合性和单位，但除明确声明外不要求交换。

\ProseParagraphBreak
包络代数与 bar 复形可对照 Weibel 的《An Introduction to Homological Algebra》\cite[Section 9.1]{Weibel1994HomologicalAlgebra}。低阶意义见\cite[Lemma 9.2.1, Proposition 9.2.2; Classification Theorem 9.3.1]{Weibel1994HomologicalAlgebra}；矩阵代数与光滑交换代数的比较分别见\cite[Section 9.5; Theorem 9.4.7]{Weibel1994HomologicalAlgebra}。杯积与括号的插入计算可对照\cite[Sections 2, 5--8]{Gerstenhaber1963CohomologyStructure}，代数形变的原始研究见\cite{Gerstenhaber1964DeformationRings}。本章从乘法展开推导形变条件，保留特征二时仍成立的障碍公式。循环同调的长正合列与混合复形分别可对照\cite[Proposition 9.6.11; Section 9.8]{Weibel1994HomologicalAlgebra}。

% !TeX root = ../../Book4.tex
% 纲要标题：### 13.1 包络代数方法
% 纲要标签：b4:sec:1201
% 纲要位置：计划纲要.md，第 3531 行。

\section{包络代数方法}
\label{b4:sec:1201}

代数乘法同时产生左、右两个作用，包络代数将它们写成一个模结构。Hochschild 理论由此既能用消解定义，也能写成显式的张量链和多线性余链。

% 纲要内容（仅作写作计划，尚非教材正文）：
%
% * 双模与包络代数
% * Hochschild 同调和上同调
% * bar 消解
%

\subsection{双模与包络代数}
\label{b4:subsec:1201:01}

一个代数的乘法同时允许从左、从右乘元素。若要用模的消解研究乘法，就应把这两种作用放在同一个模范畴内。本章始终以域 \(k\) 为底域，代数均为非零的结合含幺 \(k\) 代数，代数同态保持单位；张量积未标底环时均在 \(k\) 上进行。代数不必交换，也不必有限维。

\begin{definition}\label{b4:def:enveloping-algebra}
设 \(k\) 为域，\(A\) 为结合含幺 \(k\) 代数。其反代数 \(A^{\mathrm{op}}\) 有相同的向量空间，而乘法次序反转。代数
\[
A^e=A\otimes_k A^{\mathrm{op}},\qquad
(a\otimes b^{\mathrm{op}})(c\otimes d^{\mathrm{op}})
=ac\otimes(db)^{\mathrm{op}}
\]
称为 \(A\) 的\term[baoluo-daishu]{包络代数}[enveloping algebra]。
一个 \(k\) 向量空间 \(M\) 若带有含幺左、右 \(A\) 模作用，满足 \((am)b=a(mb)\)，且两侧的 \(k\) 作用都是原有标量乘法，便称为 \(A\) 双模。它对应的左 \(A^e\) 模作用为
\[
(a\otimes b^{\mathrm{op}})m=amb.
\]
\end{definition}
\symidx{\(A^{\mathrm{op}}\)：反代数}
\symidx{\(A^e\)：结合代数的包络代数}

这确实给出模作用，因为先用 \(c\otimes d^{\mathrm{op}}\) 作用，再用 \(a\otimes b^{\mathrm{op}}\) 作用，得到
\[
m\longmapsto cmd\longmapsto a(cmd)b=acm(db).
\]
左边的乘数按 \(ac\) 排列，右边的乘数却按 \(db\) 排列，这正是反代数出现的原因。反过来，限制到 \(A\otimes1\) 和 \(1\otimes A^{\mathrm{op}}\) 就恢复左右作用。因两个子代数彼此交换，两种作用相容。因此双模及双模同态组成的范畴就是左 \(A^e\) 模范畴。

\ProseParagraphBreak
张量积计算还需要右模。对同一个双模 \(M\)，另记 \(M^{\mathrm r}\) 为右 \(A^e\) 模，其作用规定为
\[
m\cdot(a\otimes b^{\mathrm{op}})=bma.
\]
先作用 \(a\otimes b^{\mathrm{op}}\)，再作用 \(c\otimes d^{\mathrm{op}}\)，所得为 \(dbmac\)，与乘积 \(ac\otimes(db)^{\mathrm{op}}\) 的作用一致。上标 \(\mathrm r\) 只提醒这里更换了作用方向，并未更换底层向量空间。计算 \(\operatorname{Ext}_{A^e}^n(A,M)\) 时，系数 \(M\) 保持前面的左 \(A^e\) 模结构；计算 \(\operatorname{Tor}^{A^e}_n(M^{\mathrm r},A)\) 时两个输入必须一右一左，才把第一个输入，即系数双模，改写为 \(M^{\mathrm r}\)。

\subsection{Hochschild 同调和上同调}
\label{b4:subsec:1201:02}

\begin{definition}\label{b4:def:hochschild-derived}
设 \(k\) 为域，\(A\) 为结合含幺 \(k\) 代数，\(M\) 为 \(A\) 双模，\(n\ge0\)。记 \(A^e=A\otimes_kA^{\mathrm{op}}\)，把 \(A\) 本身看作左 \(A^e\) 模，并以 \(m\cdot(a\otimes b^{\mathrm{op}})=bma\) 将 \(M\) 看作右 \(A^e\) 模 \(M^{\mathrm r}\)。定义
\[
\operatorname{HH}_n(A,M)=\operatorname{Tor}^{A^e}_n(M^{\mathrm r},A),
\qquad
\operatorname{HH}^n(A,M)=\operatorname{Ext}_{A^e}^n(A,M).
\]
它们分别称为带系数 \(M\) 的 \term[hochschild-tongdiao]{Hochschild 同调}[Hochschild homology]和 \term[hochschild-shangtongdiao]{Hochschild 上同调}[Hochschild cohomology]。当 \(M=A\) 时简记为 \(\operatorname{HH}_n(A)\)、\(\operatorname{HH}^n(A)\)。
\symidx{\(\operatorname{HH}_n(A,M),\operatorname{HH}^n(A,M)\)：Hochschild 同调与上同调}
\end{definition}

这里的 \(A\) 带有自身乘法给出的双模结构，即 \((a\otimes b^{\mathrm{op}})x=axb\)；它作为定义中 Ext、Tor 的一个固定输入，不能任意更换这一作用。零次已能看出两种不同的问题。一个双模同态 \(A\to M\) 由 \(1\) 的像 \(m\) 决定，而左右线性要求 \(am=ma\)。另一方面，\(M^{\mathrm r}\otimes_{A^e}A\) 把左右乘法产生的差消去。这两种零次函子的高次导出，分别测量它们没有保持正合的部分。

\ProseParagraphBreak
这里把底环限制为域具有实际作用：稍后 bar 消解的各项都将是 \(A^e\) 自由模。在一般交换底环上，同一个 bar 复形未必绝对投射，不能不加条件地把它计算的相对理论等同于上述 Ext、Tor。这个区别可参看 Weibel\cite[Lemma 9.1.3, Corollary 9.1.5]{Weibel1994HomologicalAlgebra}。

\subsection{bar 消解}
\label{b4:subsec:1201:03}

我们需要一个直接由乘法写出的消解，而不只是任意选择自由满射。对 \(n\ge0\)，令
\[
B_n(A)=A^{\otimes(n+2)},\qquad
\mathrm{d}_n(a_0\otimes\cdots\otimes a_{n+1})
=\sum_{i=0}^{n}(-1)^i
a_0\otimes\cdots\otimes a_i a_{i+1}\otimes\cdots\otimes a_{n+1},
\]
其中 \(\mathrm{d}_n\) 在 \(n\ge1\) 时使用；增广 \(\varepsilon:B_0(A)\to A\) 为乘法。左右 \(A\) 作用分别放在最外面的两个因子上。

\begin{proposition}\label{b4:prop:algebra-bar-resolution}
设 \(k\) 为域，\(A\) 为结合含幺 \(k\) 代数，\(A^e=A\otimes_kA^{\mathrm{op}}\)。令 \(B_n(A)=A^{\otimes_k(n+2)}\)，其中左右 \(A\) 作用分别在最外侧因子上，微分为相邻因子相乘的交错和，增广 \(B_0(A)\to A\) 为乘法。所得增广链复形 \(B_\bullet(A)\to A\) 是左 \(A^e\) 模的自由消解，称为 \(A\) 的双模 bar 消解。
\end{proposition}
\symidx{\(B_\bullet(A)\)：代数的双模 bar 消解}
\begin{proof}
最低次的微分为
\[
\mathrm{d}_1(a_0\otimes a_1\otimes a_2)
=a_0a_1\otimes a_2-a_0\otimes a_1a_2.
\]
增广后两项分别成为 \((a_0a_1)a_2\) 与 \(a_0(a_1a_2)\)，由结合律相消。一般地，两次合并不相邻因子时，两种次序得到同一张量而符号相反；合并三个连续因子时，结合律使两项相等而符号相反。因此 \(\mathrm{d}^2=0\)，增广与微分也相容。

各项有左 \(A^e\) 模同构
\[
A^e\otimes_k A^{\otimes n}\longrightarrow B_n(A),\qquad
(a\otimes b^{\mathrm{op}})\otimes(x_1\otimes\cdots\otimes x_n)
\longmapsto a\otimes x_1\otimes\cdots\otimes x_n\otimes b.
\]
逆映射把 \(a_0\otimes\cdots\otimes a_{n+1}\) 送到 \((a_0\otimes a_{n+1}^{\mathrm{op}})\otimes(a_1\otimes\cdots\otimes a_n)\)。两次复合都恢复原张量；最外两因子的左右乘法也正对应 \(A^e\) 在第一因子上的左乘，故确为左模同构。为中间向量空间选一组 \(k\) 基，便得到一组 \(A^e\) 基。

在底层向量空间上定义
\[
s_{-1}(a)=1\otimes a,\qquad
s_n(a_0\otimes\cdots\otimes a_{n+1})
=1\otimes a_0\otimes\cdots\otimes a_{n+1}.
\]
记 \(t=a_0\otimes\cdots\otimes a_{n+1}\)，并在 \(n=0\) 时以 \(\mathrm d_0=\varepsilon\) 表示增广。将前面插入的 \(1\) 与 \(a_0\) 相乘给出首项 \(t\)，其余项为
\[
\begin{aligned}
\mathrm d_{n+1}s_n(t)
&=t+\sum_{i=0}^n(-1)^{i+1}
 1\otimes a_0\otimes\cdots\otimes a_i a_{i+1}\otimes\cdots\otimes a_{n+1},\\
s_{n-1}\mathrm d_n(t)
&=\sum_{i=0}^n(-1)^i
 1\otimes a_0\otimes\cdots\otimes a_i a_{i+1}\otimes\cdots\otimes a_{n+1}.
\end{aligned}
\]
两式逐项相消，所以 \(\mathrm{d}s+s\mathrm{d}=1\)；增广处则有 \(\varepsilon s_{-1}(a)=\varepsilon(1\otimes a)=a\)。故底层向量空间复形正合，也即双模复形正合。这里的 \(s_n\) 通常不是双模同态。例如在 \(A=k[x]\) 中，\(s_{-1}(x)=1\otimes x\)，而左乘 \(x\) 后的 \(x s_{-1}(1)=x\otimes1\)，两者不同。收缩只负责证明正合性；各项的 \(A^e\) 自由性则负责保证投射性，两者合起来才说明这是可用于计算 Ext、Tor 的投射消解，并不要求 \(A\) 本身是投射双模。
\end{proof}

施加 \(\operatorname{Hom}_{A^e}(-,M)\) 得到上链复形
\[
C^n(A,M)=\operatorname{Hom}_k(A^{\otimes n},M),\qquad C^0(A,M)=M,
\]
这里将双模同态限制到 \(1\otimes a_1\otimes\cdots\otimes a_n\otimes1\)；反过来，任一 \(k\) 线性映射 \(\phi:A^{\otimes n}\to M\) 唯一延拓为 \(\widetilde\phi(a_0\otimes\cdots\otimes a_{n+1})=a_0\cdot\phi(a_1,\ldots,a_n)\cdot a_{n+1}\)。这个公式保持两侧作用，两次限制、延拓互逆。微分写成
\begin{equation}\label{b4:eq:hochschild-coboundary}
\begin{aligned}
(\delta f)(a_1,\ldots,a_{n+1})
={}&a_1\cdot f(a_2,\ldots,a_{n+1})\\
&+\sum_{i=1}^{n}(-1)^i
 f(a_1,\ldots,a_i a_{i+1},\ldots,a_{n+1})\\
&+(-1)^{n+1}f(a_1,\ldots,a_n)\cdot a_{n+1}.
\end{aligned}
\end{equation}
它计算 \(\operatorname{HH}^n(A,M)\)。特别地，对 \(k\) 线性映射 \(f:A\to M\)，
\[
(\delta f)(a,b)=a\cdot f(b)-f(ab)+f(a)\cdot b.
\]
所以 \(\delta f=0\) 正是导子恒等式。这里采用消解上通常的预合成微分，与\cref{b4:def:ext}一致；它不同于把消解当上链复形后直接使用内部 Hom 所带的逐次符号。
\symidx{\(C^n(A,M)\)：Hochschild 上链群}

\ProseParagraphBreak
作张量时，先将最外面的两个因子写成包络代数的作用：
\[
a_0\otimes\cdots\otimes a_{n+1}
=(a_0\otimes a_{n+1}^{\mathrm{op}})
\cdot(1\otimes a_1\otimes\cdots\otimes a_n\otimes1).
\]
由张量积的平衡关系和 \(M^{\mathrm r}\) 的右作用，便有
\[
\begin{aligned}
m\otimes_{A^e}(a_0\otimes\cdots\otimes a_{n+1})
&=(a_{n+1}ma_0)\otimes_{A^e}
(1\otimes a_1\otimes\cdots\otimes a_n\otimes1).
\end{aligned}
\]
因此得到同构
\[
\begin{aligned}
M^{\mathrm r}\otimes_{A^e}B_n(A)&\longrightarrow M\otimes_k A^{\otimes n},\\
m\otimes_{A^e}(a_0\otimes\cdots\otimes a_{n+1})
&\longmapsto a_{n+1}ma_0\otimes a_1\otimes\cdots\otimes a_n.
\end{aligned}
\]
其逆把 \(m\otimes a_1\otimes\cdots\otimes a_n\) 送到
\(m\otimes_{A^e}(1\otimes a_1\otimes\cdots\otimes a_n\otimes1)\)。任取 \(c\otimes d^{\mathrm{op}}\in A^e\)，平衡关系两端在此识别下的系数同为 \(a_{n+1}dmca_0\)，所以该映射良定义。

\ProseParagraphBreak
将 bar 微分施加于 \(1\otimes a_1\otimes\cdots\otimes a_n\otimes1\)，首项把 \(a_1\) 移到最左侧，经上述识别成为 \(ma_1\)；末项把 \(a_n\) 移到最右侧，因而成为 \(a_nm\)。于是 \(C_n(A,M)=M\otimes_k A^{\otimes n}\) 上的微分为
\begin{equation}\label{b4:eq:hochschild-boundary}
\begin{aligned}
b(m\otimes a_1\otimes\cdots\otimes a_n)
={}&ma_1\otimes a_2\otimes\cdots\otimes a_n\\
&+\sum_{i=1}^{n-1}(-1)^i
m\otimes\cdots\otimes a_i a_{i+1}\otimes\cdots\otimes a_n\\
&+(-1)^n a_n m\otimes a_1\otimes\cdots\otimes a_{n-1}.
\end{aligned}
\end{equation}
它计算 \(\operatorname{HH}_n(A,M)\)。特别地，
\[
\begin{aligned}
b(m\otimes a)&=ma-am,\\
b(m\otimes a\otimes c)&=ma\otimes c-m\otimes ac+cm\otimes a.
\end{aligned}
\]
二次边界的末项来自 \(m\cdot(1\otimes c^{\mathrm{op}})=cm\)，不能改成 \(mc\)。
\symidx{\(C_n(A,M)\)：Hochschild 链群}

% !TeX root = ../../Book4.tex
% 纲要标题：### 13.2 低阶 Hochschild 同调、上同调与形变
% 纲要标签：b4:sec:1202
% 纲要位置：计划纲要.md，第 3537 行。

\section{低阶 Hochschild 同调、上同调与形变}
\label{b4:sec:1202}

零次与一次 Hochschild 群记录中心、交换子商及导子；二次群则记录乘法的一阶改变。把结合律展开到一次项，才能确定这些形变何时等价。

% 纲要内容（仅作写作计划，尚非教材正文）：
%
% * 中心、导子与内导子
% * 一阶乘法形变
% * HH^2 的分类意义及等价关系
%

\subsection{中心、导子与内导子}
\label{b4:subsec:1202:01}

\begin{proposition}\label{b4:prop:hochschild-low-degrees}
设 \(k\) 为域，\(A\) 为结合含幺 \(k\) 代数，\(M\) 为 \(A\) 双模。令 \([A,M]\) 为所有 \(am-ma\)、\(a\in A\)、\(m\in M\) 的 \(k\) 线性张成空间，则
\[
\operatorname{HH}_0(A,M)=M/[A,M],\qquad
\operatorname{HH}^0(A,M)=\{\,m\in M\mid am=ma\;\text{对所有}\;a\in A\,\}.
\]
此外，令 \(\operatorname{Der}_k(A,M)\) 为满足
\(D(ab)=aD(b)+D(a)b\) 的 \(k\) 线性映射组成的空间，令
\(\operatorname{Inn}_k(A,M)\) 为形如 \(a\mapsto am-ma\) 的映射组成的子空间，则
\[
\operatorname{HH}^1(A,M)
=\operatorname{Der}_k(A,M)/\operatorname{Inn}_k(A,M).
\]
这里的映射分别称为\term[daozi]{导子}[derivation]和\term[neidaozi]{内导子}[inner derivation]。
\end{proposition}
\symidx{\(\operatorname{Der}_k(A,M),\ \operatorname{Inn}_k(A,M)\)：取值于双模的导子空间与内导子子空间}
\symidx{\([A,M]\)：双模中交换子的线性张成空间}
\begin{proof}
零次同调由 \(b(m\otimes a)=ma-am\) 得到；零次上同调是微分 \(\delta:M\to C^1(A,M)\) 的核，其中 \((\delta m)(a)=am-ma\)，所以取的是使所有这些值为零的 \(m\in M\)。由\eqref{b4:eq:hochschild-coboundary}，\(\delta D=0\) 等价于导子恒等式，而一次边界正是内导子。将这两个空间取商即得最后一个等式。
\end{proof}

取 \(M=A\)，零次上同调就是中心 \(Z(A)\)。它要求左右乘法相等；零次同调则把左右乘法的差取商，两种操作并不相同。\Cref{b4:prop:upper-triangular-low-hh}将给出中心一维而零次同调二维的代数。注意 \([A,A]\) 在同调公式中只是交换子的线性张成，通常不是理想，因此 \(\operatorname{HH}_0(A)\) 首先是向量空间商，不能径直把它当作交换化代数。
\symidx{\(Z(A)\)：结合代数的中心}

\begin{example}\label{b4:exm:polynomial-derivations}
设 \(A=k[x]\)，系数双模为 \(A\) 本身。任一导子由 \(D(x)=f(x)\) 唯一决定：Leibniz 公式归纳给出 \(D(x^r)=rx^{r-1}f(x)\)，故 \(D=f(x)\cdot\dfrac{\mathrm{d}}{\mathrm{d}x}\)。反过来，该公式定义导子。因 \(A\) 交换，内导子全为零，故 \(\operatorname{HH}^1(A)=A\cdot\dfrac{\mathrm{d}}{\mathrm{d}x}\)。此结论在任意特征成立；正特征下 \(D(x^p)=0\) 并不意味着 \(D=0\)。
\end{example}

一次同调在交换情形也有熟悉的解释。对交换 \(k\) 代数 \(A\)，用符号 \(\mathrm{d}a\) 生成 \(A\) 模，并施加 \(k\) 线性关系及
\(\mathrm{d}(ab)=a\cdot\mathrm{d}b+b\cdot\mathrm{d}a\)，所得模记为 \(\Omega_{A/k}\)，称为 \term[kahler-weifen-mo]{Kähler 微分模}[module of Kähler differentials]。这一定义直接使 \(k\) 导子 \(A\to N\) 与 \(A\) 线性映射 \(\Omega_{A/k}\to N\) 一一对应。
\symidx{\(\Omega_{A/k}\)：Kähler 微分模}

\begin{proposition}\label{b4:prop:hh1-differentials}
设 \(k\) 为域，\(A\) 为交换含幺 \(k\) 代数。映射 \(a\otimes b\mapsto a\cdot\mathrm{d}b\) 诱导自然同构
\(\operatorname{HH}_1(A)\cong\Omega_{A/k}\)。
\end{proposition}
\begin{proof}
一次微分为 \(b(a\otimes b)=ab-ba\)，交换性使它为零，所以每个一次张量都能代表同调类。二次边界是
\(ab\otimes c-a\otimes bc+ca\otimes b\)，其像因 Leibniz 关系为零。为构造逆映射，先考虑 \(D(b)=[1\otimes b]\)。它关于 \(b\) 是 \(k\) 线性的；同调的左 \(A\) 模作用在张量第一因子上，所以 \([a\otimes b]=a[1\otimes b]\)。而
\[
b(1\otimes a\otimes b)
=a\otimes b-1\otimes ab+b\otimes a
\]
给出 \([1\otimes ab]=a[1\otimes b]+b[1\otimes a]\)，即所需的 Leibniz 关系。又由 \(b(1\otimes1\otimes1)=1\otimes1\) 得 \([1\otimes1]=0\)，故 \([1\otimes c]=0\) 对所有 \(c\in k\) 成立。因此 \(D\) 是 \(k\) 导子。由微分模的泛性质，\(\mathrm db\mapsto D(b)\) 唯一延拓为 \(A\) 线性映射 \(\Omega_{A/k}\to\operatorname{HH}_1(A)\)；两映射在生成元上互逆。
\end{proof}

\subsection{一阶乘法形变}
\label{b4:subsec:1202:02}

保留一个向量空间，只允许乘法发生很小的变化，最简单的做法是引入满足 \(t^2=0\) 的参数。参数二次项全部消失以后，结合律成为关于未知映射的线性方程。

\begin{definition}\label{b4:def:first-order-deformation}
设 \(k\) 为域，\(A\) 为非零结合含幺 \(k\) 代数，\(D=k[t]/(t^2)\)。在自由 \(D\) 模 \(A\otimes_kD\) 上给定 \(D\) 双线性乘法 \(*\)。若此乘法结合、单位是原来的 \(1\)，且模 \(t\) 还原为 \(A\) 的原乘法，则称 \(*\) 为 \(A\) 的一个\term[yijie-xingbian]{一阶形变}[first-order deformation]。这里始终固定底层模及其模 \(t\) 与 \(A\) 的识别。
\end{definition}

\(D\) 双线性与还原条件说明乘法唯一写成
\[
a*b=ab+t\cdot\phi(a,b),\qquad \phi\in C^2(A,A).
\]
一般元素的乘积随之为
\((a+tu)*(b+tv)=ab+t(av+ub+\phi(a,b))\)。这不是另添的选择，而是由双线性强制得到的公式。

\begin{proposition}\label{b4:prop:first-order-cocycles}
设 \(k\) 为域，\(A\) 为非零结合含幺 \(k\) 代数，\(\phi:A\otimes_kA\to A\) 为 \(k\) 线性映射。在 \(A\otimes_k k[t]/(t^2)\) 上以双线性延拓定义乘法 \(a*b=ab+t\cdot\phi(a,b)\)。它给出单位保持为 \(1\) 的一阶形变，当且仅当
\[
\delta\phi=0,\qquad \phi(1,a)=\phi(a,1)=0\quad(a\in A).
\]
后两个条件称为 \(\phi\) 的归一化条件。
\end{proposition}
\begin{proof}
先完成内层乘法，再按双线性展开外层乘法，得到
\[
\begin{aligned}
(a*b)*c&=abc+t\bigl(\phi(ab,c)+\phi(a,b)\cdot c\bigr),\\
a*(b*c)&=abc+t\bigl(a\cdot\phi(b,c)+\phi(a,bc)\bigr).
\end{aligned}
\]
其中 \(\phi(ab,c)\) 是外层乘法的修正，\(\phi(a,b)\cdot c\) 是内层修正参与原乘法的项。比较一次系数，得到
\[
\phi(ab,c)+\phi(a,b)\cdot c=a\cdot\phi(b,c)+\phi(a,bc),
\]
正是 \(\delta\phi=0\)。常数项因原乘法结合而相等，更高项因 \(t^2=0\) 消失。单位条件则等价于 \(1*a=a*1=a\)，即所列两个等式。
\end{proof}

上述形变允许新的乘法不交换，即使 \(A\) 原本交换。若还要求形变交换，就必须再加 \(\phi(a,b)=\phi(b,a)\)；不能把所有 Hochschild 二次类都解释为交换代数的形变。

\subsection{\texorpdfstring{\(\operatorname{HH}^2\) 与形变的分类}{HH2 与形变的分类}}
\label{b4:subsec:1202:03}

一个看起来非零的一次改动，有时只是把原来的元素换了一套坐标。因而分类之前必须说清允许怎样更换坐标。

\begin{definition}\label{b4:def:deformation-equivalence}
设 \(k\) 为域，\(A\) 为非零结合含幺 \(k\) 代数，\(D=k[t]/(t^2)\)，\(*_\phi,*_\psi\) 是 \(A\) 的一阶形变。若存在保持单位的 \(D\) 线性代数同构
\[
T:(A\otimes_kD,*_\phi)\longrightarrow(A\otimes_kD,*_\psi),
\qquad T\bmod t=1_A,
\]
则称两个形变等价。这样的映射唯一写成 \(T=1+t h\)，其中 \(h:A\to A\) 为 \(k\) 线性映射且 \(h(1)=0\)；其逆为 \(1-th\)。
\end{definition}

\begin{theorem}\label{b4:thm:hh2-deformation-classification}
设 \(k\) 为域，\(A\) 为非零结合含幺 \(k\) 代数。保持单位的一阶形变在\cref{b4:def:deformation-equivalence}的等价关系下，与 \(\operatorname{HH}^2(A)\) 典范一一对应。零类对应原乘法的平凡形变。
\end{theorem}
\begin{proof}
两种乘法都在模 \(t\) 后成为原乘法，\(T=1+th\) 则只改变一阶坐标。把源乘法搬到目标乘法时，必须使先相乘再换坐标与先换坐标再相乘一致。展开 \(T(a*_\phi b)=T(a)*_\psi T(b)\)，一次系数满足
\[
\phi(a,b)+h(ab)=\psi(a,b)+a\cdot h(b)+h(a)\cdot b.
\]
故等价恰好要求 \(\phi-\psi=\delta h\)，并且 \(h(1)=0\)。

还需证明归一化并未丢失二次上同调类。给定任意\(2\)-余圈 \(\phi\)，在 \(\delta\phi=0\) 中代入 \((1,1,a)\)、\((a,1,1)\)，得到
\[
0=\phi(1,a)-\phi(1,1)a,\qquad
0=a\phi(1,1)-\phi(a,1).
\]
两式中的 \(\phi(1,1)\) 是同一个元素，记作 \(c\)，便有
\[
\phi(1,a)=ca,\qquad \phi(a,1)=ac,\qquad c=\phi(1,1).
\]
因 \(k\) 是域且 \(1\ne0\)，可选线性映射 \(h:A\to A\) 使 \(h(1)=c\)。余边公式给出 \(\delta h(1,a)=h(1)a=ca\)、\(\delta h(a,1)=ah(1)=ac\)，所以 \(\phi-\delta h\) 在 \((1,a)\)、\((a,1)\) 上都为零，即归一化，并代表同一上同调类。若两个归一化余圈之差为 \(\delta h\)，在 \((1,1)\) 上求值便得 \(h(1)=0\)，恰好符合允许的坐标变换。因而余圈、上同调类和形变等价类之间的对应既满又单。
\end{proof}

这一定理分类的是带固定还原识别的形变。若允许 \(T\bmod t\) 是 \(A\) 的任意自同构，还要进一步考虑 \(\operatorname{Aut}_k(A)\) 对 \(\operatorname{HH}^2(A)\) 的作用。
\symidx{\(\operatorname{Aut}_k(A)\)：\(A\) 的 \(k\) 代数自同构群}

\ProseParagraphBreak
更一般的系数双模也有二次上同调的解释。设 \(k\) 为域，\(A\) 为非零结合含幺 \(k\) 代数，\(M\) 为指定的 \(A\) 双模。考虑含幺 \(k\) 代数的平方零扩张
\[
0\longrightarrow M\xrightarrow{i}E\xrightarrow{q}A\longrightarrow0,
\qquad i(M)^2=0,
\]
其中固定核、商与 \(M,A\) 的识别，并要求 \(E\) 的乘法在核上诱导指定的双模作用：若 \(q(e)=a\)，则 \(e\,i(m)=i(am)\)、\(i(m)e=i(ma)\)。平方零条件使这些作用与提升 \(e\) 的选择无关。两个扩张等价，是指存在保持这两端识别的代数同构。

\ProseParagraphBreak
因 \(k\) 是域，先将 \(1\) 扩充为 \(A\) 的一组 \(k\) 基，再为其余基向量选取原像，就得到满足 \(s(1)=1\) 的 \(k\) 线性截面 \(s:A\to E\)。这只分裂了向量空间；\(s\) 未必保持乘法，乘法误差正是要记录的数据。用 \(s(a)+i(m)\) 把底层向量空间识别为 \(A\oplus M\)，乘法成为
\[
(a,m)(b,n)=(ab,\,an+mb+\phi(a,b)),
\qquad i(\phi(a,b))=s(a)s(b)-s(ab).
\]
结合律恰好给出 \(\delta\phi=0\)，单位条件给出 \(\phi(1,a)=\phi(a,1)=0\)。反过来，每个这样的归一化 \(2\)-余圈都用此公式给出一个扩张。将截面改成 \(s'=s+i\circ h\)，其中 \(h(1)=0\)，因 \(i(M)^2=0\)，新的误差为
\[
\phi'(a,b)=\phi(a,b)+a\cdot h(b)+h(a)\cdot b-h(ab)
=\phi(a,b)+(\delta h)(a,b).
\]
同一个元素在旧坐标中是 \((a,m)\)，在新坐标中是 \((a,m-h(a))\)。因此这里的新旧截面给出的加号，与形变分类中从源到目标的坐标变换 \(\phi-\psi=\delta h\) 相容。与形变分类证明相同的归一化计算适用于系数 \(M\)，故这些扩张的等价类与 \(\operatorname{HH}^2(A,M)\) 典范一一对应。零类对应分裂扩张 \(A\oplus M\)。这个一般系数版本亦见 Weibel\cite[Classification Theorem 9.3.1]{Weibel1994HomologicalAlgebra}。

\begin{example}\label{b4:exm:dual-numbers-first-deformation}
设 \(\operatorname{char}k\ne2\)，\(A=k[x]/(x^2)\)。归一化二次上链只需指定 \(\phi(x,x)=u+vx\)，且它自动为余圈：唯一需检验的三元组是 \((x,x,x)\)，在此 \(\delta\phi=x\cdot\phi(x,x)-\phi(x,x)\cdot x=0\)。归一化一上链满足 \(h(1)=0\)，而
\(\delta h(x,x)=2x\cdot h(x)\) 遍历 \(kx\)。所以 \(\operatorname{HH}^2(A)\cong k\)，可由 \(\phi(x,x)=1\) 代表。对应乘法为 \(x*x=t\)，即 \(D[x]/(x^2-t)\)，其 \(D\) 基仍为 \(1,x\)。
\end{example}

\begin{proposition}\label{b4:prop:upper-triangular-low-hh}
设 \(k\) 为任意域，\(T\) 为 \(k\) 上的二阶上三角矩阵代数。则
\[
\operatorname{HH}^0(T)=kI_2,\qquad
\operatorname{HH}_0(T)\cong k^2,
\]
其中第二个同构由两个对角元给出。
\end{proposition}
\begin{proof}
设 \(z\in Z(T)\)。与 \(E_{11}\) 交换迫使 \(z\) 的非对角项为零，再与 \(E_{12}\) 交换迫使两个对角元相等，故 \(Z(T)=kI_2\)。两个上三角矩阵的交换子对角项总为零，而 \([E_{11},E_{12}]=E_{12}\)，故 \([T,T]=kE_{12}\)。由\cref{b4:prop:hochschild-low-degrees}，两个结论随即成立。这些计算不需除以 \(2\)，在特征二同样有效。
\end{proof}

% !TeX root = ../../Book4.tex
% 纲要标题：### 13.3 乘积与括号
% 纲要标签：b4:sec:1203
% 纲要位置：计划纲要.md，第 3543 行。

\section{乘积与括号}
\label{b4:sec:1203}

余链值相乘与将余链代入另一个余链，是两种次数不同的运算。后者把结合律的高阶项写成括号，使延拓障碍成为具体的三次上同调类。

% 纲要内容（仅作写作计划，尚非教材正文）：
%
% * 杯积
% * Gerstenhaber 括号的构造思想
% * 形变障碍与 HH^3 的作用
%

\subsection{杯积}
\label{b4:subsec:1203:01}

当系数是 \(A\) 本身时，可以把两个上链的值相乘。另一种做法是把一个上链的值代入另一个上链。前者给出杯积，后者适合表达结合律及形变障碍；两种运算的次数不同。

\begin{definition}\label{b4:def:hochschild-cup}
设 \(k\) 为域，\(A\) 为结合含幺 \(k\) 代数，\(f\in C^p(A,A)\)、\(g\in C^q(A,A)\)，其中 \(p,q\ge0\)。规定
\[
(f\smile g)(a_1,\ldots,a_{p+q})
=f(a_1,\ldots,a_p)g(a_{p+1},\ldots,a_{p+q}).
\]
这个运算称为\term[beiji]{杯积}[cup product]。当某次为零时，相应上链就是一个 \(A\) 中元素，占据该因子的位置。
\end{definition}
\symidx{\(f\smile g\)：Hochschild 上链及上同调类的杯积}

\begin{proposition}\label{b4:prop:hochschild-cup-differential}
设 \(k\) 为域，\(A\) 为结合含幺 \(k\) 代数，\(p,q\ge0\)，\(f\in C^p(A,A)\)、\(g\in C^q(A,A)\)。\Cref{b4:def:hochschild-cup}定义的杯积结合，零次上链 \(1\) 是单位，且
\[
\delta(f\smile g)=(\delta f)\smile g+(-1)^p f\smile\delta g.
\]
因而它诱导 \(\operatorname{HH}^*(A)\) 上的带单位分次结合代数结构。
\end{proposition}
\begin{proof}
结合性来自 \(A\) 的乘法结合性。先看 \(p=q=1\) 的微分：
\[
\begin{aligned}
\delta(f\smile g)(a,b,c)
={}&a\cdot f(b)g(c)-f(ab)g(c)\\
&+f(a)g(bc)-f(a)g(b)\cdot c\\
={}&(\delta f)(a,b)g(c)-f(a)(\delta g)(b,c).
\end{aligned}
\]
右边展开时，\(f(a)b g(c)\) 的两次出现符号相反；这是一般次数中交界项相消的最小情形。在微分公式中，合并前 \(p+1\) 个位置产生 \((\delta f)\smile g\) 的各内部项，合并后 \(q+1\) 个位置产生 \((-1)^pf\smile\delta g\) 的各内部项。两部分交界处额外出现的项分别为
\((-1)^{p+1}f(a_1,\ldots,a_p)a_{p+1}\cdot g(\ldots)\) 与
\((-1)^pf(a_1,\ldots,a_p)a_{p+1}\cdot g(\ldots)\)，恰好抵消；最左、最右作用项保留下来，即为 \(\delta(f\smile g)\)。零次情形由同一左右乘法公式成立。

因此两个余圈的杯积仍是余圈。若 \(f=\delta u\)、\(\delta g=0\)，则 \(f\smile g=\delta(u\smile g)\)；另一因子为余边时同样成立。这证明乘法在上同调类上良定义。
\end{proof}

杯积也是 \(\operatorname{Ext}_{A^e}^*(A,A)\) 的 Yoneda 乘积。为直接比较两种定义，设 \(g\in C^q(A,A)\) 为余圈，在 bar 消解上定义
\[
\begin{aligned}
\widetilde g_i:B_{q+i}(A)&\longrightarrow B_i(A),\\
a_0\otimes\cdots\otimes a_{q+i+1}
&\longmapsto a_0\otimes\cdots\otimes a_i\otimes
g(a_{i+1},\ldots,a_{i+q})a_{q+i+1}\qquad(i\ge0).
\end{aligned}
\]
这是双模同态；\(\delta g=0\) 使合并尾部因子所得各项相消，因而 \(\mathrm{d}_i\widetilde g_i=\widetilde g_{i-1}\mathrm{d}_{q+i}\) 对 \(i\ge1\) 成立。增广与 \(\widetilde g_0\) 的复合就是代表 \(g\) 的 bar 余圈。若 \(f\in C^p(A,A)\) 也是余圈，把它与 \(\widetilde g_p\) 复合，在标准张量上便得到 \(f(a_1,\ldots,a_p)g(a_{p+1},\ldots,a_{p+q})\)。这正是\cref{b4:prop:yoneda-resolution-product}中用消解提升计算 Yoneda 乘积的公式。

\ProseParagraphBreak
例如中心中的元素通过零次杯积作用于所有上同调群。一般系数双模 \(M\) 没有自带的 \(M\otimes M\to M\)，所以不能不加数据就为 \(\operatorname{HH}^*(A,M)\) 写出同样的乘法。

\subsection{Gerstenhaber 括号的构造思想}
\label{b4:subsec:1203:02}

杯积把两组自变量前后排列，不能直接表达 \(\phi(\phi(a,b),c)\)。为比较这类嵌套乘法，需要记录插入的位置及其符号。例如两个二次上链的插入将是
\[
(f\circ g)(a,b,c)=f(g(a,b),c)-f(a,g(b,c)).
\]
一次插入把两个自变量合成一个值，所以结果有三个自变量。一般地，插入 \(q\) 次上链会给后续位置增加 \(q-1\) 个自变量；符号所使用的正是这个移后次数。

\begin{definition}\label{b4:def:gerstenhaber-bracket}
设 \(k\) 为域，\(A\) 为结合含幺 \(k\) 代数，\(p,q\ge0\)。对 \(f\in C^p(A,A)\)、\(g\in C^q(A,A)\)，在 \(p\ge1\) 时定义
\[
\begin{aligned}
f\circ g&=\sum_{i=0}^{p-1}(-1)^{i(q-1)}f\circ_i g,\\
(f\circ_i g)(a_1,\ldots,a_{p+q-1})
&=f(a_1,\ldots,a_i,g(a_{i+1},\ldots,a_{i+q}),
a_{i+q+1},\ldots,a_{p+q-1}).
\end{aligned}
\]
\(q=0\) 时将常数 \(g\in A\) 插入第 \(i+1\) 个位置；\(p=0\) 时规定 \(f\circ g=0\)。定义
\[
[f,g]=f\circ g-(-1)^{(p-1)(q-1)}g\circ f.
\]
这个次数为 \(-1\) 的运算称为 \term[gerstenhaber-kuohao]{Gerstenhaber 括号}[Gerstenhaber bracket]。\footnote{格斯滕哈伯（Murray Gerstenhaber，1927-2024），美国数学家，研究代数形变理论，并引入以他命名的括号与代数结构。}两个零次上链的括号取零。
\end{definition}
\symidx{\([f,g]\)：Hochschild 上链上的 Gerstenhaber 括号}

零次上链也参与插入。例如 \(D\in C^1(A,A)\)、\(a\in C^0(A,A)=A\) 满足 \(D\circ a=D(a)\)、\(a\circ D=0\)，故 \([D,a]=D(a)\)。它是一个零次上链，次数仍为 \(1+0-1\)。

\begin{proposition}\label{b4:prop:gerstenhaber-identities}
设 \(k\) 为域，\(A\) 为结合含幺 \(k\) 代数，\(p,q\ge0\)，以 \(|f|=p-1\) 表示 \(f\in C^p(A,A)\) 的移后次数。\Cref{b4:def:gerstenhaber-bracket}定义的括号满足分次反对称性及分次 Jacobi 恒等式，且诱导
\[
\operatorname{HH}^p(A)\times\operatorname{HH}^q(A)
\longrightarrow\operatorname{HH}^{p+q-1}(A).
\]
这里约定负次数的上同调为零。若 \(m(a,b)=ab\) 为原乘法，则
\[
[m,f]=(-1)^{p-1}\delta f.
\]
此外，杯积在上同调上满足分次交换律
\[
\alpha\smile\beta=(-1)^{pq}\beta\smile\alpha,
\qquad \alpha\in\operatorname{HH}^p(A),\quad
\beta\in\operatorname{HH}^q(A).
\]
\end{proposition}
\begin{proof}
记 \(a=|f|\)、\(b=|g|\)、\(c=|h|\)，并令
\[
P(f,g,h)=(f\circ g)\circ h-f\circ(g\circ h).
\]
先令 \(f,g,h\) 都为二次上链，观察两种实际复合。左边的树把 \(h\) 插入 \(g\) 内部，再把 \(g\) 插入 \(f\)；右边把 \(g,h\) 插入 \(f\) 的两个不同位置。箭头表示把下方运算的输出代入上方运算。
\[
\begin{tikzpicture}
\node (nf) at (-3,0) {\(f\)};
\node (ng) at (-3.7,-0.9) {\(g\)};
\node (na4) at (-2.2,-0.9) {\(a_4\)};
\node (nh) at (-4.2,-1.8) {\(h\)};
\node (na3) at (-3.1,-1.8) {\(a_3\)};
\node (na1) at (-4.7,-2.7) {\(a_1\)};
\node (na2) at (-3.7,-2.7) {\(a_2\)};
\draw[->] (ng)--(nf); \draw[->] (na4)--(nf);
\draw[->] (nh)--(ng); \draw[->] (na3)--(ng);
\draw[->] (na1)--(nh); \draw[->] (na2)--(nh);
\node (df) at (2.3,0) {\(f\)};
\node (dg) at (1.4,-0.9) {\(g\)};
\node (dh) at (3.2,-0.9) {\(h\)};
\node (da1) at (0.9,-1.8) {\(a_1\)};
\node (da2) at (1.9,-1.8) {\(a_2\)};
\node (da3) at (2.7,-1.8) {\(a_3\)};
\node (da4) at (3.7,-1.8) {\(a_4\)};
\draw[->] (dg)--(df); \draw[->] (dh)--(df);
\draw[->] (da1)--(dg); \draw[->] (da2)--(dg);
\draw[->] (da3)--(dh); \draw[->] (da4)--(dh);
\end{tikzpicture}
\]
左边的 \(f(g(h(a_1,a_2),a_3),a_4)\) 同时出现在 \((f\circ g)\circ h\) 与 \(f\circ(g\circ h)\) 中，符号均为正，所以在 \(P\) 中抵消；把 \(h\) 插入 \(g\) 的另一个位置时，两处符号均为负，也抵消。右边的 \(f(g(a_1,a_2),h(a_3,a_4))\) 则不属于 \(f\circ(g\circ h)\)。在 \((f\circ g)\circ h\) 中，先占 \(f\) 的位置 \(0\) 后，\(h\) 应插入新位置 \(2\)，符号为 \((-1)^0(-1)^2=1\)；若先在 \(f\) 的位置 \(1\) 插入 \(h\)，再在位置 \(0\) 插入 \(g\)，符号为 \((-1)^1(-1)^0=-1\)。因此这一次展开后只剩
\[
\begin{aligned}
P(f,g,h)(a_1,a_2,a_3,a_4)
={}&f(g(a_1,a_2),h(a_3,a_4))\\
&-f(h(a_1,a_2),g(a_3,a_4)).
\end{aligned}
\]
交换 \(g,h\) 恰使这个表达式变号。一般次数的计算整理的是同样的嵌套抵消与异位配对。插入后的自变量数为 \(p+q-1\)，减去一后为 \((p-1)+(q-1)\)，所以移后次数恰使插入成为次数相加的运算。

以下 \(i,j\) 表示从零开始计数的插入位置，而 \(a,b,c\) 表示移后次数；它们与图中的自变量 \(a_1,\ldots,a_4\) 无关。先在 \(f\) 的位置 \(i\) 插入 \(g\)，再在其中 \(g\) 的位置 \(j\) 插入 \(h\)，符号是下式左端；先在 \(g\) 中插入 \(h\)，再把所得余链插入 \(f\) 的位置 \(i\)，符号是下式右端：
\[
(-1)^{ib+(i+j)c}
=(-1)^{i(b+c)+jc}.
\]
因此这类内部插入项在 \(P(f,g,h)\) 中抵消，只剩把 \(g,h\) 放入 \(f\) 的不同位置的项。

若 \(g\) 占原位置 \(i\)，\(h\) 占原位置 \(j\)，且 \(i<j\)，先插 \(g\) 后插 \(h\) 时，\(h\) 的位置变为 \(j+b\)，符号为
\[
(-1)^{ib+(j+b)c};
\]
先插 \(h\) 后插 \(g\) 的符号为 \((-1)^{jc+ib}\)，两者之比是 \((-1)^{bc}\)。当 \(j<i\) 时，两种符号分别为
\((-1)^{ib+jc}\) 与 \((-1)^{jc+(i+c)b}\)，仍相差 \((-1)^{bc}\)。于是
\[
P(f,g,h)=(-1)^{bc}P(f,h,g).
\]
这就是带符号的右 pre-Lie 恒等式。若 \(g\) 或 \(h\) 为零次上链，内部位置为空，插入常数则使后续位置前移一格，即相应移后次数为 \(-1\)；上面的异位配对仍成立。若 \(f\) 为零次上链，两边均为零。所有落入负上链次数的表达式也取零。

括号的分次反对称性直接来自定义。为核验 Jacobi 恒等式，把其导子形式的差展开为
\[
\begin{aligned}
&[f,[g,h]]-[[f,g],h]-(-1)^{ab}[g,[f,h]]\\
={}&-P(f,g,h)+(-1)^{bc}P(f,h,g)\\
&+(-1)^{ab}\bigl(P(g,f,h)-(-1)^{ac}P(g,h,f)\bigr)\\
&-(-1)^{c(a+b)}\bigl(P(h,f,g)-(-1)^{ab}P(h,g,f)\bigr).
\end{aligned}
\]
第一对由 \(P(f,g,h)=(-1)^{bc}P(f,h,g)\) 抵消；后两对分别交换 \(f,h\) 和 \(f,g\)，由同一恒等式抵消，故差为零。再用分次反对称性改写，得到通常的循环形式
\[
(-1)^{|f|\cdot|h|}[f,[g,h]]
+(-1)^{|g|\cdot|f|}[g,[h,f]]
+(-1)^{|h|\cdot|g|}[h,[f,g]]=0.
\]
将二次上链 \(m\) 代入括号定义，最外两项是左右乘法，内部各项是合并相邻自变量，给出 \([m,f]=(-1)^{p-1}\delta f\)。例如对二次上链 \(\phi\)，
\[
\begin{aligned}
[m,\phi](a,b,c)
={}&\phi(a,b)\cdot c-a\cdot\phi(b,c)\\
&+\phi(ab,c)-\phi(a,bc)
=-(\delta\phi)(a,b,c).
\end{aligned}
\]
原乘法的结合律就是 \(m\circ m=0\)，而一阶结合律就是 \([m,\phi]=0\)。特别地，令 \(D f=[m,f]\)，Jacobi 恒等式\footnote{雅可比（Carl Gustav Jacob Jacobi，1804-1851），德国数学家，在椭圆函数和动力学的微分方程方面作出重要贡献。}给出
\[
D[f,g]=[Df,g]+(-1)^{|f|}[f,Dg].
\]
\(D\) 与 \(\delta\) 在每次只差一个符号，因此有相同的余圈、余边。最后一个等式说明余圈的括号仍为余圈，且一个因子为余边、另一个为余圈时结果为余边。这就使括号下降到上同调。

最后把原乘法放在异位插入公式的首位。由于 \(m\) 只有两个位置，其余两项给出
\[
P(m,f,g)=(-1)^{p(q-1)}f\smile g+(-1)^{p-1}g\smile f.
\]
展开 \([m,f\circ g]\)，由 \(P\) 的定义改写括号位置，再用
\(P(f,m,g)=(-1)^{q-1}P(f,g,m)\) 抵消这两项，得到
\[
[m,f\circ g]=[m,f]\circ g+(-1)^{p-1}f\circ[m,g]-P(m,f,g).
\]
这里 \(f\circ g\) 的次数为 \(p+q-1\)，故左端是 \((-1)^{p+q-2}\delta(f\circ g)\)。右端分别代入 \([m,f]=(-1)^{p-1}\delta f\)、\([m,g]=(-1)^{q-1}\delta g\) 及前式的两个杯积，再将两边乘以 \((-1)^{p+q-2}\)，得到
\[
\begin{aligned}
\delta(f\circ g)
={}&(-1)^{q-1}(\delta f)\circ g+f\circ\delta g\\
&+(-1)^{q-1}\bigl((-1)^{pq}f\smile g-g\smile f\bigr).
\end{aligned}
\]
若 \(f,g\) 都是余圈，右边最后一行就是余边，故上同调上的杯积满足分次交换律。例如对导子 \(D,E\)，这个等式成为
\[
\delta(D\circ E)=-(D\smile E+E\smile D).
\]
两个杯积在上链层面未必互为负号，其和却是明确的余边，所以上同调类互为负号。零次情形同样包含在一般公式中；当两者都为零次余圈时，它们是中心元素，杯积直接交换。
\end{proof}

杯积与括号还满足分次 Leibniz 律：对
\(\alpha\in\operatorname{HH}^p(A)\)、\(\beta\in\operatorname{HH}^q(A)\)、\(\gamma\in\operatorname{HH}^r(A)\)，有
\[
[\alpha\smile\beta,\gamma]
=[\alpha,\gamma]\smile\beta
+(-1)^{p(r-1)}\alpha\smile[\beta,\gamma].
\]
按本节的括号符号约定，固定右变量 \(\gamma\) 得到杯积的次数为 \(r-1\) 的分次导子。这使 \(\operatorname{HH}^*(A)\) 同时具有相容的分次交换乘法与次数为 \(-1\) 的 Lie 括号，合称 Gerstenhaber 代数结构。\cite[Section 8, Theorem 5, Corollaries 1--2, pp. 285--287]{Gerstenhaber1963CohomologyStructure}
\begin{proofsketch}
取代表 \(\alpha,\beta,\gamma\) 的余圈 \(f,g,h\)。将 \(h\) 插入杯积的前半或后半，直接得到
\[
(f\smile g)\circ h
=(f\circ h)\smile g+(-1)^{p(r-1)}f\smile(g\circ h).
\]
另一方向需比较 \(h\circ(f\smile g)\) 与
\((-1)^{q(r-1)}(h\circ f)\smile g+f\smile(h\circ g)\)。所用的双重插入上链可以写成
\[
H=\sum_{i=0}^{r-2}\ \sum_{j=i+p}^{p+r-2}
(-1)^{(p-1)i+(q-1)j}(h\circ_i f)\circ_j g.
\]
先把 \(f\) 插入 \(h\) 的位置 \(i\)，再把 \(g\) 插入其右方的另一个原位置；前次插入增加了 \(p-1\) 个位置，故第二个指标从 \(i+p\) 开始。\(h\) 的最后一个原位置是 \(r-1\)，移动后成为 \((r-1)+(p-1)=p+r-2\)，这给出上界；为了给右方的 \(g\) 留一个原位置，\(i\) 只取到 \(r-2\)。两次插入各带 \((-1)^{(p-1)i}\)、\((-1)^{(q-1)j}\)，相乘即为和式中的符号。\(H\) 的次数为 \(p+q+r-2\)，在 \(r<2\) 时取空和。展开 \(\delta H\) 时，内部合并项由 \(\delta f=\delta g=\delta h=0\) 抵消，异位项按位置配对。剩下的端项分别对应把两个相邻插入合成杯积、让 \(g\) 留在最右端、让 \(f\) 留在最左端；归并后的等式为
\[
\begin{aligned}
\delta H=(-1)^{q(p-1)}\bigl(&h\circ(f\smile g)
 -(-1)^{q(r-1)}(h\circ f)\smile g\\
 &-f\smile(h\circ g)\bigr).
\end{aligned}
\]
因此所比较的差恰为 \((-1)^{q(p-1)}\delta H\)，是一条余边。结合杯积插入等式与括号定义，便得到所需 Leibniz 律。这里省略 \(\delta H\) 的全部位置符号展开；完整计算见上述文献。
\end{proofsketch}

对一次上链，即线性自映射，括号为普通交换子 \([D,E]=D\circ E-E\circ D\)。因此导子的交换子仍是导子，并给出 \(\operatorname{HH}^1(A)\) 上的 Lie 括号\footnote{李（Marius Sophus Lie，1842-1899），挪威数学家，创立连续变换群理论，奠定 Lie 群与 Lie 代数的基础。}。对二次上链 \(\phi\)，则
\[
(\phi\circ\phi)(a,b,c)
=\phi(\phi(a,b),c)-\phi(a,\phi(b,c)).
\]
当 \(\operatorname{char}k\ne2\) 时，它等于 \(\dfrac12[\phi,\phi]\)。在特征二中不能用除以二的写法；实际障碍始终是左边这个插入表达式。

\subsection{\texorpdfstring{形变障碍与 \(\operatorname{HH}^3\)}{形变障碍与 HH3}}
\label{b4:subsec:1203:03}

一次结合律只给线性方程，第二次就出现刚才的二次表达式。所有阶的形变需要先确定形式级数的含义。设 \(k\) 为域，\(A\) 为结合含幺 \(k\) 代数，记
\[
A[\![t]\!]
=\varprojlim_{N\ge0}\bigl(A\otimes_k k[t]/(t^{N+1})\bigr)
\cong\prod_{n\ge0}At^n.
\]
其元素为 \(\sum_{n\ge0}a_nt^n\)，\(t\) 进拓扑以 \(t^{N+1}A[\![t]\!]\) 为零点邻域基。若 \(A\) 无限维，这一完备化一般不同于普通张量积 \(A\otimes_k k[\![t]\!]\)。

\ProseParagraphBreak

例如取 \(A=k[x]\)，则 \(\sum_{n\ge0}x^nt^n\) 是 \(A[\![t]\!]\) 中的元素。普通张量积中的元素却只能是有限和 \(\sum_{i=1}^r a_i\otimes f_i(t)\)。把它展开为 \(t\) 的级数后，每个系数都落在同一个有限维空间 \(\operatorname{span}_k\{a_1,\ldots,a_r\}\) 中；而 \(1,x,x^2,\ldots\) 线性无关，不可能全部落在这种空间里。因此该级数不来自普通张量积，完备化允许的是这样的无限系数族。

\ProseParagraphBreak
在这一完备模上，保持单位的形式形变是连续的 \(k[\![t]\!]\) 双线性结合乘法，模 \(t\) 后等于 \(A\) 的原乘法，单位仍为 \(1\)。它在常数元素上的值唯一写成
\[
a*b=ab+\sum_{i\ge1}t^i\mu_i(a,b),
\]
其中 \(\mu_i\) 为归一化 \(k\) 双线性映射。两个这样的乘法若由保持单位、模 \(t\) 为恒等的连续 \(k[\![t]\!]\) 线性同构 \(T\) 相联系，即
\(T(a*b)=T(a)*'T(b)\)，就称为等价。这样的坐标变换写成
\(T=\operatorname{id}+\sum_{i\ge1}t^ih_i\)，其中 \(h_i:A\to A\) 为 \(k\) 线性映射，且 \(h_i(1)=0\)；其逆可在各有限截断中逐次求出。

\ProseParagraphBreak
令 \(\mu_0=m\) 为原乘法。双线性与连续性使一般级数的乘积必须为
\[
\left(\sum_{r\ge0}a_rt^r\right)*
\left(\sum_{s\ge0}b_st^s\right)
=\sum_{n\ge0}t^n
\sum_{\substack{i+r+s=n\\i,r,s\ge0}}\mu_i(a_r,b_s).
\]
固定 \(n\) 时，内层只取满足 \(i+r+s=n\) 的有限多个三元组。因此每个系数都已确定，模 \(t^{N+1}\) 的乘积只依赖两个输入的相应截断；这也说明该公式给出连续的 \(k[\![t]\!]\) 双线性乘法。

\ProseParagraphBreak
若已知 \(\phi\) 是归一化 \(2\)-余圈，尝试令
\[
a*b=ab+t\cdot\phi(a,b)+t^2\cdot\psi(a,b)\pmod {t^3}.
\]
比较结合律的二次系数，得到
\[
\begin{aligned}
\psi(ab,c)+\psi(a,b)\cdot c+\phi(\phi(a,b),c)
={}&\psi(a,bc)+a\cdot\psi(b,c)\\
&+\phi(a,\phi(b,c)).
\end{aligned}
\]
前四个 \(\psi\) 项来自二阶修正与原乘法的结合，两个嵌套 \(\phi\) 项则来自两次一阶修正。将 \(\psi\) 项移到一侧，所需且仅需的条件是
\begin{equation}\label{b4:eq:second-deformation-obstruction}
\delta\psi=\phi\circ\phi.
\end{equation}
因此出现一个三次上同调问题：右边不仅要是余圈，还必须是余边。

\begin{theorem}\label{b4:thm:hochschild-obstruction}
设 \(k\) 为域，\(A\) 为非零结合含幺 \(k\) 代数，\(m(a,b)=ab\) 为其乘法，\(N\ge1\)，\(D_N=k[t]/(t^{N+1})\)。在 \(A\otimes_kD_N\) 上给定保持单位的结合 \(D_N\) 双线性乘法
\[
m_t=m+t\mu_1+\cdots+t^N\mu_N,
\]
其中各 \(\mu_i:A\otimes_kA\to A\) 为 \(k\) 线性映射，满足 \(\mu_i(1,a)=\mu_i(a,1)=0\) 对所有 \(a\in A\) 成立。令
\[
O_{N+1}(a,b,c)=
\sum_{\substack{i+j=N+1\\i,j\ge1}}
\bigl(\mu_i(\mu_j(a,b),c)-\mu_i(a,\mu_j(b,c))\bigr).
\]
则 \(O_{N+1}\) 是归一化三余圈。该给定形变能延拓到模 \(t^{N+2}\) 的保持单位形变，当且仅当
\([O_{N+1}]=0\) 于 \(\operatorname{HH}^3(A)\)。
\end{theorem}
\begin{proof}
在现有截断多项式乘法后加上 \(t^{N+1}\nu\)，结合子的 \(t^{N+1}\) 系数为 \(O_{N+1}-\delta\nu\)，因为涉及原乘法 \(m\) 与 \(\nu\) 的四项恰好是 \(-\delta\nu\)。故延拓要求 \(\delta\nu=O_{N+1}\)。

为证明右边是余圈，不需要除以任何整数。对任意双线性乘法 \(*\)，记其结合子为 \(F(a,b,c)=(a*b)*c-a*(b*c)\)。直接展开可得
\[
\begin{aligned}
0={}&a*F(b,c,d)-F(a*b,c,d)+F(a,b*c,d)\\
&-F(a,b,c*d)+F(a,b,c)*d.
\end{aligned}
\]
展开后出现的五种四重乘积是
\[
\begin{gathered}
((a*b)*c)*d,\qquad (a*(b*c))*d,\qquad (a*b)*(c*d),\\
a*((b*c)*d),\qquad a*(b*(c*d)).
\end{gathered}
\]
每一种恰出现两次，符号相反；这一步未使用 \(*\) 的结合性。对尚未加入 \(\nu\) 的 \(m_t\) 使用此式。因结合子的所有低于 \(N+1\) 次系数为零，比较 \(t^{N+1}\) 系数时，外部乘法只取原乘法，所得正是 \(\delta O_{N+1}=0\)。各 \(\mu_i\) 归一化也使 \(O_{N+1}\) 任一自变量为 \(1\) 时为零。

若延拓存在，\(O_{N+1}=\delta\nu\) 已说明其类为零。反过来，若类为零，先取任意 \(\nu\) 使 \(\delta\nu=O_{N+1}\)。右边在带单位的三元组上为零，所以与\cref{b4:thm:hh2-deformation-classification}中相同的两次代入给出
\(\nu(1,a)=ca\)、\(\nu(a,1)=ac\)，其中 \(c=\nu(1,1)\)。选择 \(h(1)=c\)，以 \(\nu-\delta h\) 替代 \(\nu\)，即可既保持 \(\delta\nu=O_{N+1}\)，又使 \(\nu\) 归一化。所得延拓保持单位。
\end{proof}

\begin{corollary}\label{b4:cor:hh3-unobstructed}
设 \(k\) 为域，\(A\) 为非零结合含幺 \(k\) 代数，\(m(a,b)=ab\) 为其乘法。若 \(\operatorname{HH}^3(A)=0\)，则每个保持单位的一阶形变都能逐次延拓成 \(A[\![t]\!]\) 上的结合形式乘法
\(m_t=m+\sum_{i\ge1}t^i\mu_i\)，其中各 \(\mu_i:A\otimes_kA\to A\) 为归一化 \(k\) 线性映射。
\end{corollary}
\begin{proof}
从给定的一阶形变 \(m^{(1)}\) 开始，用\cref{b4:thm:hochschild-obstruction}归纳选择 \(\mu_{N+1}\)，得到保持单位的结合截断乘法 \(m^{(N+1)}\)，并使其模 \(t^{N+1}\) 恰为已选的 \(m^{(N)}\)。因此这些乘法构成兼容系统，可取逆极限。所得级数乘法就是前面的一般系数公式；它连续且 \(k[\![t]\!]\) 双线性。每个次数的结合律已在相应截断中成立，单位由所有 \(\mu_i\) 的归一化保证。
\end{proof}

这是形式延拓，不涉及实数或复数分析中的收敛。若 \(\operatorname{HH}^3(A)\ne0\)，也不能反向断言所有形变都会受阻；要计算的是当前已选系数产生的特定类 \([O_{N+1}]\)。后续系数的不同选择还可能影响更高阶障碍。下面的代数说明，三次上同调非零也可能与所有一阶形变都可延拓同时发生。

\begin{proposition}\label{b4:prop:dual-numbers-unobstructed}
设 \(k\) 为任意域，\(A=k[x]/(x^2)\)。则 \(\operatorname{HH}^3(A)\ne0\)，但 \(A\) 的每个保持单位的一阶形变都能延拓成 \(A[\![t]\!]\) 上的保持单位形式形变。
\end{proposition}
\begin{proof}
定义三线性上链 \(\eta\)，使 \(\eta(x,x,x)=x\)，且在基 \(1,x\) 组成的三元组中，只要有一个自变量为 \(1\)，其值就为零。归一化使 \(\delta\eta\) 在带单位的四元组上为零；在唯一剩下的基四元组上，
\[
(\delta\eta)(x,x,x,x)
=x\eta(x,x,x)+\eta(x,x,x)x=2x^2=0.
\]
因此 \(\eta\) 是余圈。另一方面，对任意二次上链 \(\psi\)，由 \(x^2=0\) 和交换性得
\[
(\delta\psi)(x,x,x)=x\psi(x,x)-\psi(x,x)x=0.
\]
它不可能等于取值为 \(x\ne0\) 的 \(\eta\)，故 \([\eta]\ne0\)。

任一保持单位的一阶形变由归一化余圈 \(\phi(x,x)=u+vx\)、\(u,v\in k\) 确定。取结合代数
\[
k[\![t]\!][X]/\bigl(X^2-t(u+vX)\bigr).
\]
首一多项式除法使它在 \(k[\![t]\!]\) 上以 \(1,X\) 为自由基，故可与 \(A[\![t]\!]\) 识别；其乘法连续，单位保持不变，模 \(t^2\) 后恰为给定的一阶形变。这在任何特征下都成立。
\end{proof}

\ProseParagraphBreak
低阶群的含义可以合在一起比较：零次记录中心化与交换子商，一次记录导子与微分，二次分类一阶乘法形变，三次容纳逐次延拓的障碍类。具体对象、系数和条件见\cref{b4:tab:hochschild-low-degrees}。
\begin{table}[htbp]
	\centering
	\caption[Hochschild 理论的低次数解释]{Hochschild 理论的低次数解释。这里 \(k\) 为域，\(A\) 为非零结合含幺 \(k\) 代数，\(M\) 为 \(A\) 双模。}
	\label{b4:tab:hochschild-low-degrees}
	\begin{tabularx}{\textwidth}{@{}p{0.22\textwidth}>{\raggedright\arraybackslash}X@{}}
		\toprule
		对象 & 代数意义 \\
		\midrule
		\(\operatorname{HH}_0(A,M)\) & \(M/[A,M]\)，消去左右作用之差 \\
		\(\operatorname{HH}^0(A,M)\) & \(\{\,m\in M\mid am=ma\;\text{对所有}\;a\in A\,\}\) \\
		\(\operatorname{HH}_1(A)\) & 当 \(A\) 交换时，等于 Kähler 微分模 \(\Omega_{A/k}\) \\
		\(\operatorname{HH}^1(A,M)\) & \(\operatorname{Der}_k(A,M)/\operatorname{Inn}_k(A,M)\) \\
		\(\operatorname{HH}^2(A)\) & 保持单位及固定还原识别的一阶乘法形变等价类 \\
		\(\operatorname{HH}^2(A,M)\) & 固定核 \(M\)、商 \(A\) 及双模作用的平方零代数扩张等价类 \\
		\(\operatorname{HH}^3(A)\) & 给定有限阶形变的延拓障碍类所在的空间；并非每个三次类都必须出现为障碍 \\
		\bottomrule
	\end{tabularx}
\end{table}


% !TeX root = ../../Book4.tex
% 纲要标题：### 13.4 典型计算、Morita 不变性与 HKR
% 纲要标签：b4:sec:1204
% 纲要位置：计划纲要.md，第 3549 行。

\section{典型计算、Morita 不变性与 HKR}
\label{b4:sec:1204}

矩阵代数、光滑交换代数及一个平方零代数展示三种不同的 Hochschild 计算。前两种分别利用模范畴等价和对角的 Koszul 消解，最后一种直接检验形变障碍。交换代数的形变还会产生一个新的问题：乘法最先出现的不交换部分，满足哪些恒等式？Poisson 括号正是从低阶结合律中留下的结构。

% 纲要内容（仅作写作计划，尚非教材正文）：
%
% * 矩阵代数与 Morita 不变性
% * 光滑交换代数的比较方向
% * 不把一阶可形变误判为一定可延拓至所有阶
%
% * 新增小节状态：初稿已完成；原有内容保留。
%
% #### 13.4.4 Poisson 括号与形式星积
%
% ---
%

\subsection{矩阵代数与 Morita 不变性}
\label{b4:subsec:1204:01}

矩阵代数的乘法看起来比标量代数复杂，但模范畴之间的等价会保留相应的同调信息。下面用矩阵单位写出这一等价，并从零次交换子商的比较得到 Hochschild 同调的 Morita 不变性\footnote{\OriginalHanName{森田紀一}（假名：もりた きいち，罗马音：Kiiti Morita，1915-1995），日本数学家，建立以模范畴等价为核心的森田理论。}。

\begin{theorem}\label{b4:thm:hh-matrix-morita}
	设 \(k\) 为域，\(A\) 为结合含幺 \(k\) 代数，\(r\ge1\)，\(B=M_r(A)\)。对任意 \(A\) 双模 \(N\)，令 \(M_r(N)\) 带有左右矩阵乘法作用。则对所有 \(n\ge0\)，有自然同构
	\[
		\operatorname{HH}_n(B,M_r(N))\cong\operatorname{HH}_n(A,N),\qquad \operatorname{HH}^n(B,M_r(N))\cong\operatorname{HH}^n(A,N).
	\]
\end{theorem}
\begin{proof}
	记 \(E_{ij}\) 为矩阵单位，\(e=E_{11}\)，并识别 \(eBe=A\)。对 \(M=M_r(N)\)，\(E_{1i}mE_{j1}\) 正是将 \(m\) 的第 \((i,j)\) 个分量放进 \((1,1)\) 角部，所以该角部足以恢复整个矩阵双模。对一般 \(B\) 双模，仍可用相同的矩阵单位来恢复各分量。双模函子
	\[
		F(M)=eMe,\qquad G(N)=M_r(N)
	\]
	互为等价。具体地，\(FG(N)=N\)，而
	\[
		M\longrightarrow M_r(eMe),\qquad m\longmapsto(E_{1i}mE_{j1})_{i,j}
	\]
	的逆将 \((n_{ij})\) 送到 \(\sum_{i,j}E_{i1}n_{ij}E_{1j}\)。矩阵单位等式 \(E_{ij}E_{uv}=\delta_{ju}E_{iv}\) 验证两次复合为恒等，且这些映射保持双模作用。两个函子均正合：若 \(f:M\twoheadrightarrow M'\) 为双模满射，给定 \(em'e\)，先取 \(f(m)=m'\)，再取角部 \(eme\)，就有 \(f(eme)=em'e\)；核则满足 \(\ker F(f)=e(\ker f)e\)：若 \(x\in eMe\) 且 \(f(x)=0\)，便有 \(x=exe\in e(\ker f)e\)；反向包含由 \(f(eye)=ef(y)e=0\) 得到。而 \(M_r(-)\) 是有限直和，也保持正合。它们还保持投射对象，因为一个提升问题经逆等价转移后仍是满射的提升问题。

	对 \(B\) 取投射双模消解 \(P_\bullet\to B\)。施加 \(F\) 后得到 \(A\) 的投射双模消解，而且等价的全忠实性给出复形同构
	\[
		\operatorname{Hom}_{B^e}(P_\bullet,M) \cong\operatorname{Hom}_{A^e}(F(P_\bullet),F(M)).
	\]
	取上同调便证明上同调断言。

	上同调刚才使用的是代数本身的消解；同调则可以借\cref{b4:thm:tor-balanced}改用系数双模的消解。这样便能逐项使用双模范畴等价。为此先考察右正合函子 \(Q_B(M)=M/[B,M]\)。它与张量函子的联系由以下互逆映射给出：
	\[
		\begin{aligned}
			M^{\mathrm r}\otimes_{B^e}B&\longrightarrow Q_B(M), &m\otimes b&\longmapsto[bm],\\
			Q_B(M)&\longrightarrow M^{\mathrm r}\otimes_{B^e}B, &[m]&\longmapsto m\otimes1.
		\end{aligned}
	\]
	第一式满足平衡关系：取 \(\gamma=a\otimes c^{\mathrm{op}}\in B^e\)，则 \(m\cdot\gamma=cma\)、\(\gamma\cdot b=abc\)，所以 \((cma)\otimes b=m\otimes abc\)。两者的像为 \([bcma]\)、\([abcm]\)；在交换子商中可将最右端的 \(a\) 移到最左端，因为它们的差 \((bcm)a-a(bcm)\) 是一个交换子。第二式消去交换子，因为
	\[
		\begin{aligned}
			bm\otimes1 &=m\cdot(1\otimes b^{\mathrm{op}})\otimes1 =m\otimes b,\\
			mb\otimes1 &=m\cdot(b\otimes1)\otimes1 =m\otimes b.
		\end{aligned}
	\]
	此外 \(bm\otimes1=m\otimes b\) 验证两次复合为恒等。因此 \(Q_B(M)\cong M^{\mathrm r}\otimes_{B^e}B\)。从双模到右 \(B^e\) 模的函子 \(M\mapsto M^{\mathrm r}\) 是正合等价，保持投射对象，所以在系数双模 \(M\) 的投射消解上计算，并用\cref{b4:thm:tor-balanced}，得到
	\[
		L_nQ_B(M) \cong\operatorname{Tor}^{B^e}_n(M^{\mathrm r},B) =\operatorname{HH}_n(B,M).
	\]

	现在比较两个交换子商。矩阵单位给出自然同构
	\[
		Q_B(M)\longrightarrow Q_A(eMe),\qquad [m]\longmapsto\left[\sum_iE_{1i}mE_{i1}\right].
	\]
	其良定义可直接检验。将 \(1=\sum_jE_{j1}E_{1j}\) 插入 \(b\) 与 \(m\) 之间，得到
	\[
		E_{1i}bmE_{i1}=\sum_j(E_{1i}bE_{j1})(E_{1j}mE_{i1}).
	\]
	这两个因子分别属于 \(eBe=A\) 与 \(eMe\)，所以在目标交换子商中可以交换先后。交换后再对两个指标求和，由 \(E_{i1}E_{1i}=E_{ii}\) 及 \(\sum_iE_{ii}=1_B\) 得到
	\[
		\begin{aligned}
			\sum_{i,j}(E_{1j}mE_{i1})(E_{1i}bE_{j1}) &=\sum_jE_{1j}m\Bigl(\sum_iE_{ii}\Bigr)bE_{j1}\\
			&=\sum_jE_{1j}mbE_{j1}.
		\end{aligned}
	\]
	这就说明 \([bm]\) 与 \([mb]\) 有相同的像。逆映射由角部的包含 \(eMe\subseteq M\) 给出。在 \(Q_B(M)\) 中，关系
	\[
		E_{ii}m-E_{1i}mE_{i1}=[E_{i1},E_{1i}m]
	\]
	求和后说明逆向复合是恒等。对 \(x\in eMe\)，有 \(x=exe\)；当 \(i\ne1\) 时，\(E_{1i}e=eE_{i1}=0\)，只有 \(i=1\) 的项留下，故 \(\sum_iE_{1i}xE_{i1}=x\)。正向复合也为恒等。

	所得零次比较就是 \(Q_B\cong Q_A\circ F\)。\(F\) 正合且保持投射，将 \(M\) 的双模投射消解送到 \(F(M)\) 的投射消解；逐项应用零次自然同构，再取同调，便有
	\[
		\operatorname{HH}_n(B,M) \cong L_nQ_B(M) \cong L_nQ_A(F(M)) \cong\operatorname{HH}_n(A,F(M)).
	\]
	取 \(M=M_r(N)\) 即得结论。
\end{proof}

在 \(M=M_r(N)\) 时，刚才的零次比较把矩阵 \((m_{ij})\) 的类送到 \([\sum_i m_{ii}]\)，即由矩阵迹给出；矩阵单位 \(E_{1i}mE_{i1}\) 将每个对角分量都移到同一角部。这里的迹取值于 \(N/[A,N]\)，其对循环换序的不变性来自取商，未要求 \(A\) 交换。

\ProseParagraphBreak
一般的 \(k\) 线性 Morita 等价由互逆的张量双模实现，同样保持相应系数的 Hochschild 群。矩阵情形的具体等价与\cite[Proposition 9.5.2, Corollary 9.5.3]{Weibel1994HomologicalAlgebra}一致；下面说明一般情形的零次比较为何也能导出。

\begin{proofsketch}
	设 \({}_AP_B\)、\({}_BQ_A\) 实现第二卷定理16.3.3及推论16.3.4的 Morita 等价，取相容的双模同构 \(P\otimes_BQ\cong A\)、\(Q\otimes_AP\cong B\)，分别把配对写作 \(\langle p,q\rangle_A\) 与 \(\langle q,p\rangle_B\)。双模函子
	\[
		F(N)=Q\otimes_AN\otimes_AP
	\]
	及其逆是正合等价，送 \(A\) 到 \(B\)，并保持投射对象。这里选择从 \(A\) 双模到 \(B\) 双模的方向；它与刚才从矩阵双模取角部的方向相反，而在矩阵情形对应 \(N\mapsto M_r(N)\)。因此对 \(A\) 的双模投射消解施加 \(F\)，再用等价的 Hom 双射，得到 \(\operatorname{HH}^n(A,N)\cong\operatorname{HH}^n(B,F(N))\)。

	同调端的零次比较为
	\[
		F(N)/[B,F(N)]\longrightarrow N/[A,N],\qquad [q\otimes n\otimes p]\longmapsto [n\langle p,q\rangle_A].
	\]
	张量平衡关系在右边分别变成等式或交换子关系；\(B\) 交换子关系由 \(\langle pb,q\rangle_A=\langle p,bq\rangle_A\) 消去。取有限表达 \(1_A=\sum_i\langle p_i,q_i\rangle_A\)。这是单位在 \(P\otimes_BQ\) 中的原像，而普通张量积的每个元素都是有限个纯张量之和，所以这里的有限性不要求 \(A\) 有限维。逆映射送 \([n]\) 到 \(\sum_i[q_i\otimes n\otimes p_i]\)。一向复合直接为恒等，另一向使用 Morita 配对的相容式
	\[
		\langle p,q\rangle_Ap_i=p\langle q,p_i\rangle_B, \qquad \langle q,p_i\rangle_Bq_i=q\langle p_i,q_i\rangle_A,
	\]
	以及交换子商中允许把右侧的 \(B\) 因子移到左侧，便恢复 \([q\otimes n\otimes p]\)。这也使逆与所选单位表达无关。

	上述张量识别对任意结合代数都成立，因而
	\[
		L_nQ_A(N)\cong\operatorname{HH}_n(A,N), \qquad L_nQ_B(F(N))\cong\operatorname{HH}_n(B,F(N)).
	\]
	在 \(N\) 的双模投射消解上逐项应用 \(Q_B\circ F\cong Q_A\)，再利用 \(F\) 正合且保持投射，就得到两端 Hochschild 同调的同构。这里省略的是把等价的单位、余单位写成上述相容配对的标准 Morita 构造。完整同调比较也见\cite[Theorem 9.5.6, pp. 328--329]{Weibel1994HomologicalAlgebra}，该处用双复形比较实现相同的系数变换。
\end{proofsketch}

\begin{example}\label{b4:exm:matrix-hochschild-vanishing}
	对任意域 \(k\) 和 \(r\ge1\)，\(k^e=k\)，所以
	\[
		\begin{gathered}
			\operatorname{HH}_0(M_r(k))=k,\qquad \operatorname{HH}^0(M_r(k))=k,\\
			\operatorname{HH}_n(M_r(k))=\operatorname{HH}^n(M_r(k))=0\quad(n>0).
		\end{gathered}
	\]
	零次同调的同构是普通迹在交换子商上诱导的映射，其逆送 \(c\) 到 \([cE_{11}]\)，不需要除以 \(r\)。即使特征整除 \(r\)，\([I_r]\) 的迹为零，也不妨碍整个同调空间是一维。由\cref{b4:prop:hochschild-low-degrees}，矩阵代数的导子都是内导子；由\cref{b4:thm:hh2-deformation-classification}，其一阶形变均等价于平凡形变。
\end{example}

\subsection{光滑交换代数的 Hochschild 同调}
\label{b4:subsec:1204:02}

\(\operatorname{HH}_1(A)=\Omega_{A/k}\) 提示：在合适的交换代数上，高次同调应当与微分的外幂有关。先看一个可以只用已学消解算完的情形。

\begin{proposition}\label{b4:prop:polynomial-hochschild}
	设 \(A=k[x_1,\ldots,x_d]\)，其中 \(k\) 为域，\(d\ge0\)。对每个 \(n\ge0\)，作为 \(A\) 模有
	\[
		\operatorname{HH}_n(A)\cong\bigwedge_A^n A^d, \qquad \operatorname{HH}^n(A)\cong \bigwedge_A^n\operatorname{Hom}_A(A^d,A).
	\]
	在同调端可把标准基与 \(\mathrm{d}x_1,\ldots,\mathrm{d}x_d\) 对应；两者在 \(n>d\) 时都为零。
\end{proposition}
\begin{proof}
	包络代数为 \(S=k[x_1,\ldots,x_d,y_1,\ldots,y_d]\)，\(A\) 是商 \(S/(x_1-y_1,\ldots,x_d-y_d)\)。这里 \(x_i\) 来自张量积的第一份 \(A\)，\(y_i\) 来自第二份；乘法映射把两份变量送到同一个 \(x_i\)，其核因而由这些差生成。依次令 \(x_i=y_i\) 后仍为多项式环，所以下一个 \(x_j-y_j\) 不为零因子。这是正则列，由\cref{b4:prop:koszul-resolution}，相应 Koszul 复形是 \(A\) 的有限自由 \(S\) 模消解。

	例如 \(d=2\) 时，记 \(z_i=x_i-y_i\)，该消解为
	\[
		0\longrightarrow S\xrightarrow{\mathrm{d}_2}S^2 \xrightarrow{\mathrm{d}_1}S\longrightarrow A\longrightarrow0, \quad \mathrm{d}_2(c)=(-z_2c,z_1c),\quad \mathrm{d}_1(u,v)=z_1u+z_2v.
	\]
	与 \(A\) 张量后，\(z_1,z_2\) 均为零，所以微分消失，同调依次为 \(A,A^2,A\)，对应于 \(1\)、\(\mathrm{d}x_1,\mathrm{d}x_2\) 及 \(\mathrm{d}x_1\wedge\mathrm{d}x_2\)。一般 \(d\) 时机制相同：张量后各项为 \(\bigwedge_A^n A^d\)，取 \(\operatorname{Hom}_S(-,A)\) 后则为其 \(A\) 线性对偶，两种复形的微分都为零。有限自由模的楔积基将后者识别为 \(\bigwedge_A^n\operatorname{Hom}_A(A^d,A)\)。微分模的基 \(\mathrm{d}x_i\) 则由多项式逐变量求导的泛性质得到。
\end{proof}

对一般交换代数，比较同调与微分外幂所需的条件是光滑性。它由平方零扩张中的提升性质与有限展示共同刻画。

\begin{definition}\label{b4:def:smooth-algebra-hh}
	设 \(k\) 为域，\(A\) 为有限展示的交换含幺 \(k\) 代数。若对每个交换含幺 \(k\) 代数 \(E\)、每个平方零理想 \(J\subseteq E\) 及每个 \(k\) 代数同态 \(A\to E/J\)，都存在提升 \(A\to E\)，则称 \(A\) 在 \(k\) 上\term[guanghua-daishu]{光滑}[smooth]。平方零指 \(J^2=0\)，提升不要求唯一。
\end{definition}

这个约定与第三卷定义14.5.1、14.5.2相同：平方零提升的存在性称为形式光滑，有限展示加形式光滑称为光滑。在形变问题中，用提升性质表述光滑，能够直接看出方程中的误差是否可消去。例如 \(A=k[x]/(x^3)\) 不光滑：取 \(E=k[z]/(z^4)\)、\(J=(z^3)\)，则 \(J^2=0\)，且 \(E/J\cong A\)。恒等映射若能提升，\(x\) 的像必为 \(z+c z^3\)，但其立方是 \(z^3\ne0\)，无法保持 \(x^3=0\)。这在任意特征下都给出一个具体的提升失败。

\ProseParagraphBreak
多项式代数满足此条件，因为分别提升各生成元的像即可。这里要求 \(E\) 交换；它不等同于对所有非交换扩张的提升性质，也不蕴含 \(\operatorname{HH}^2(A)=0\)。例如两个变量的多项式代数在\cref{b4:prop:polynomial-hochschild}中有非零二次上同调。

\ProseParagraphBreak
以下 Hochschild--Kostant--Rosenberg 定理\footnote{科斯坦特（Bertram Kostant，1928-2017），美国数学家，研究 Lie 群与表示论；罗森伯格（Alex F. T. W. Rosenberg，1926-2007），德裔美国数学家，研究环论和同调代数。}给出证明概要。Weibel\cite[Theorem 9.4.7]{Weibel1994HomologicalAlgebra}给出本质有限型版本；下面采用有限展示的情形。

\begin{theorem}[Hochschild--Kostant--Rosenberg 定理]\label{b4:thm:hkr-external}
	设 \(k\) 为域，\(A\) 为在 \(k\) 上光滑的有限展示交换代数。则有自然 \(A\) 模同构
	\[
		\bigwedge_A^n\Omega_{A/k}\longrightarrow\operatorname{HH}_n(A) \qquad(n\ge0).
	\]
	它把 \(a_0\cdot\mathrm{d}a_1\wedge\cdots\wedge \mathrm{d}a_n\) 送到
	\[
		\left[\sum_{\sigma\in S_n}\operatorname{sgn}(\sigma) a_0\otimes a_{\sigma(1)}\otimes\cdots\otimes a_{\sigma(n)}\right].
	\]
\end{theorem}

\begin{proofsketch}
	先构造定理中的全局映射，再局部检验它是否为同构。交换代数的 Hochschild 链带有 shuffle 乘法：将两条链的首因子相乘，把其余两串因子按各自原有次序交错排列，并赋予交错置换的符号。该乘法与边界相容，因而下降到同调；这里相乘的是同调链的类，不能使用前面将余链值相乘的杯积公式。最低阶计算为
	\[
		[1\otimes a]\,[1\otimes b] =[1\otimes a\otimes b-1\otimes b\otimes a].
	\]
	这两项分别对应 \(a,b\) 的两种交错次序。特别地，同一个一次链的 shuffle 平方为零，包括特征二的情形；这是交错项在链上配对抵消的结果，不能仅由分次交换律推出。\Cref{b4:prop:hh1-differentials}已使 \(\mathrm{d}a\mapsto[1\otimes a]\) 保持微分模的关系。shuffle 乘法因此诱导自然 \(A\) 线性映射 \(\Psi:\bigwedge_A^*\Omega_{A/k}\to\operatorname{HH}_*(A)\)。反复相乘一次类，便得到各张量因子的带符号置换和，正是定理的公式。这同时证明了公式对微分模关系和外幂关系良定义，没有使用 \(n!\) 可逆。

	写 \(A^e=A\otimes_kA\)，以乘法映射的核 \(J\) 定义对角理想。令 \(\Delta a=a\otimes1-1\otimes a\)，则
	\[
		\Delta(ab)=(a\otimes1)\Delta b+(1\otimes b)\Delta a.
	\]
	模 \(J^2\) 后，两份 \(b\) 的作用相同，故 \(a\mapsto\Delta a+J^2\) 是 \(k\) 导子，诱导 \(\Omega_{A/k}\to J/J^2\)。其逆把 \(\sum_i a_i\otimes b_i\in J\) 的类送到 \(-\sum_i a_i\cdot\mathrm{d}b_i\)；乘积法则和 \(\sum_i a_ib_i=0\) 保证它零化 \(J^2\)。两次复合在生成元 \(\mathrm{d}a\) 与 \(\Delta a+J^2\) 上为恒等，因而 \(J/J^2\cong\Omega_{A/k}\)。

	现在局部计算 \(\operatorname{HH}_n(A)=\operatorname{Tor}^{A^e}_n(A,A)\)。这里使用光滑代数的一项几何结果：在对角上的每个点，\(J\) 在相应的局部环中由正则列生成，其长度等于 \(\Omega_{A/k}\) 的局部秩。这就是光滑有限型态射的对角为正则嵌入的代数表述。从平方零提升条件得到这一局部正则列的证明，见 Weibel\cite[Proposition 9.4.6 and proof of Theorem 9.4.7, pp. 322--323]{Weibel1994HomologicalAlgebra}。

	用该正则列的 Koszul 复形消解相应的局部商，再与商张量。正则列作用为零，所以微分全部为零，局部 Tor 代数是 \(J/J^2\) 的外代数。这与多项式计算相同。Koszul 外积和 bar 的 shuffle 乘法相比较时，一次项是刚才的 \(J/J^2\cong\Omega_{A/k}\) 识别，故局部同构就是 \(\Psi\) 的局部化。两个消解的外积比较见\cite[Application 8.7.12, Example 8.7.13, pp. 291--292; Corollary 9.4.4 and proof of Theorem 9.4.7, pp. 322--323]{Weibel1994HomologicalAlgebra}。于是 \(\Psi\) 的核、余核在每个素理想处局部化都为零，因而全局也为零。概要省略的是正则对角的局部证明及 Koszul、bar 外积比较的链映射构造；全局映射本身已由 shuffle 定义，不需另选局部同构来黏合。
\end{proofsketch}

该同构不要求特征零。若另有 \(\operatorname{char}k=0\)，逆映射由下式给出：
\[
	[a_0\otimes\cdots\otimes a_n]\longmapsto \dfrac{1}{n!}a_0\cdot\mathrm{d}a_1\wedge\cdots\wedge \mathrm{d}a_n.
\]
未除以 \(n!\) 的链映射 \(a_0\otimes\cdots\otimes a_n\mapsto a_0\cdot\mathrm{d}a_1\wedge\cdots\wedge\mathrm{d}a_n\) 与正向反对称映射复合后，得到 \(n!\) 倍的恒等映射：每个置换的符号与外幂换序的符号相乘均为正号，共有 \(n!\) 项。因此逆映射的上述写法需要除以 \(n!\)，正向映射的同构性却不需要。例如特征二时，在 \(k[x,y]\) 的二次同调中，\([1\otimes x\otimes y-1\otimes y\otimes x]\) 仍对应非零的 \(\mathrm{d}x\wedge\mathrm{d}y\)，但未除以 \(2!\) 的反向映射把它送到零。上述结论比较的是同调群；整个上链复形的乘法与括号的相容性还需另行讨论。

\subsection{一阶形变的高阶延拓障碍}
\label{b4:subsec:1204:03}

\cref{b4:exm:dual-numbers-first-deformation}的一阶形变确能继续延拓：\(k[\![t]\!][x]/(x^2-t)\) 以 \(1,x\) 为自由基，结合性来自多项式商。但下面的同样小的代数已经有不能通过第二阶的形变。

\begin{example}\label{b4:exm:obstructed-deformation}
	设 \(k\) 为任意域，\(A=k\oplus kx\oplus ky\)，其中 \(V=kx\oplus ky\) 满足 \(V^2=0\)，\(1\) 为单位。定义归一化二次上链 \(\phi\)，在 \(V\times V\) 上规定
	\[
		\phi(x,x)=y,\qquad \phi(y,x)=y,\qquad \phi(x,y)=\phi(y,y)=0.
	\]
	当三个自变量属于 \(V\) 时，\(\delta\phi\) 的左右作用项因 \(V^2=0\) 消失，内部项也因相邻乘积为零而消失；某个自变量为 \(1\) 时，归一化使其余两项抵消。由三线性，\(\delta\phi=0\)。故 \(a*b=ab+t\cdot\phi(a,b)\) 确是一阶形变。

	\ProseParagraphBreak
	然而
	\[
		(\phi\circ\phi)(x,x,x) =\phi(y,x)-\phi(x,y)=y\ne0.
	\]
	对任意二次上链 \(\psi\)，交换性和 \(x^2=0\) 给出
	\[
		(\delta\psi)(x,x,x) =x\cdot\psi(x,x)-\psi(x,x)\cdot x=0.
	\]
	仅有障碍上链在 \((x,x,x)\) 处非零，还不足以断言障碍类非零；关键是所有余边在该处都取零值，而 \(\phi\circ\phi\) 取值为 \(y\)。因此它不可能是余边，\eqref{b4:eq:second-deformation-obstruction}无解。这个一阶形变在任何特征下都不能延拓至模 \(t^3\)。
\end{example}

此例把无法延拓的判定化为乘法表上的有限计算：先求\(2\)-余圈，再在一个特定三元组上计算障碍，最后检验所有余边在该三元组上的值都为零。障碍在这里取非零值，因此无论怎样选择其余二阶系数，都无法消去它。

\subsection{Poisson 括号与形式星积}
\label{b4:subsec:1304:04}

设原代数 \(A\) 交换，而形变乘法不再交换。因为零次项相同，两个常数元素的交换子必被形变参数整除。取它的一次系数，便得到原交换代数上的一个双线性运算；这个系数允许为零，不交换性也可能在更高次才出现。结合律对一次系数的限制分成两层：一次结合律产生乘积法则，二次结合律才迫使它满足 Jacobi 恒等式。

\begin{definition}\label{b4:def:poisson-algebra}
	设 \(k\) 为域，\(A\) 为交换含幺 \(k\) 代数。若 \(k\) 双线性运算 \(\{-,-\}:A\times A\to A\) 满足
	\[
		\begin{gathered}
			\{a,a\}=0,\qquad
			\{a,\{b,c\}\}+\{b,\{c,a\}\}+\{c,\{a,b\}\}=0,\\
			\{a,bc\}=\{a,b\}c+b\{a,c\}
		\end{gathered}
	\]
	对任意 \(a,b,c\in A\) 成立，就称它为 \(A\) 上的 \term[Poisson-kuohao]{Poisson 括号}[Poisson bracket]，称带有这个括号的 \(A\) 为 \term[Poisson-daishu]{Poisson 代数}[Poisson algebra]。\footnote{泊松（Siméon-Denis Poisson，1781-1840），法国数学家、物理学家，研究分析、概率及力学；力学中的泊松括号以他命名。}
\end{definition}
\symidx{\(\{a,b\}\)：Poisson 括号}

第一式称为交替性；对 \(a+b\) 使用它，得到 \(\{a,b\}=-\{b,a\}\)。因此第三式也给出第一变量的乘积法则，即括号在每个变量上都是导子。在特征二中，反对称性本身不能代替交替性。这里的括号没有分次：\(A\) 是普通交换代数，Jacobi 恒等式也不带移后次数的符号。

\begin{definition}\label{b4:def:formal-star-product}
	设 \(k\) 为域，\(A\) 为交换含幺 \(k\) 代数，\(\hbar\) 为形式参数。\(A[\![\hbar]\!]\) 上一个连续、保持单位的 \(k[\![\hbar]\!]\) 双线性结合乘法，若模 \(\hbar\) 后等于 \(A\) 的原乘法，就称为 \(A\) 上的\term[xingshi-xingji]{形式星积}[formal star product]。在常数元素上将它写成
	\[
		a\star b=ab+\sum_{r\ge1}\hbar^r C_r(a,b),\qquad
		C_r(1,a)=C_r(a,1)=0.
	\]
	这里每个 \(C_r:A\times A\to A\) 都是 \(k\) 双线性映射，连续性相对于 \(\hbar\)-进拓扑。
\end{definition}
\symidx{\(a\star b\)：形式星积}

这是\subsecref{b4:subsec:1203:03}中形式形变在交换代数上的情形；本节把参数改记为 \(\hbar\)，它仍是中心的形式参数，不表示一个待取极限的实数。一般级数的乘积仍按每个 \(\hbar^n\) 系数中的有限和定义。这里不附加 \(C_r\) 为微分算子的要求；下面的具体构造会给出由导数组成的系数。若星积的一次交换子满足
\[
	C_1(a,b)-C_1(b,a)=\{a,b\},
\]
就说它量化了这个 Poisson 括号。该条件规定了不交换性的一次项，并未唯一确定所有高阶系数。

\begin{proposition}\label{b4:prop:star-first-commutator-poisson}
	设 \(k\) 为任意域，\(A\) 为交换含幺 \(k\) 代数，\(C_1:A\times A\to A\) 为归一化 \(k\) 双线性映射。若
	\[
		a\star b=ab+\hbar C_1(a,b)\pmod{\hbar^2}
	\]
	经 \(k[\hbar]/(\hbar^2)\) 双线性延拓，给出 \(A\otimes_k k[\hbar]/(\hbar^2)\) 上保持单位的结合乘法，则
	\[
		\{a,b\}=C_1(a,b)-C_1(b,a)
	\]
	是交替双导子。若这个乘法还可延拓成模 \(\hbar^3\) 的保持单位结合乘法，则它是 Poisson 括号。特别地，每个形式星积都产生 Poisson 括号；在模 \(\hbar\) 为恒等的形变等价下，所得括号不变。
\end{proposition}
\begin{proof}
	双线性和交替性直接来自定义。记 \([u,v]_{\star}=u\star v-v\star u\)。任何结合乘法的交换子满足
	\[
		[a,b\star c]_{\star}
		=[a,b]_{\star}\star c+b\star[a,c]_{\star}.
	\]
	因为 \(A\) 的原乘法交换，\([a,b]_{\star}=\hbar\{a,b\}\pmod{\hbar^2}\)。左边输入 \(b\star c=bc+\hbar C_1(b,c)\) 的额外一项不影响一次系数：它与 \(a\) 的零次交换子为零。取上式的 \(\hbar\) 系数，便得
	\[
		\{a,bc\}=\{a,b\}c+b\{a,c\}.
	\]
	反对称性给出另一变量的乘积法则。

	现在假设存在二阶延拓，并写
	\[
		[a,b]_{\star}=\hbar\{a,b\}+\hbar^2 Q(a,b)\pmod{\hbar^3}.
	\]
	结合乘法的交换子满足 Jacobi 恒等式：把每个嵌套交换子展开后，六种有序三重乘积各出现两次且符号相反。因此
	\[
		[a,[b,c]_{\star}]_{\star}
		+[b,[c,a]_{\star}]_{\star}
		+[c,[a,b]_{\star}]_{\star}=0.
	\]
	其中 \([a,\hbar^2 Q(b,c)]_{\star}\) 模 \(\hbar^3\) 为零，因为原乘法交换；而
	\[
		[a,\hbar\{b,c\}]_{\star}
		=\hbar^2\{a,\{b,c\}\}\pmod{\hbar^3}.
	\]
	取二次系数就得到 Poisson 的 Jacobi 恒等式。若只在模 \(\hbar^2\) 的代数中计算，这些二次项全为零，不能据此推出任何 Jacobi 条件。这正是需要二阶延拓的原因。

	最后，设等价变换的一次项为 \(T=\operatorname{id}+\hbar H\pmod{\hbar^2}\)。由\cref{b4:thm:hh2-deformation-classification}中的比较公式，两种一次系数之差为
	\[
		(\delta H)(a,b)=a\cdot H(b)-H(ab)+H(a)\cdot b.
	\]
	因 \(A\) 交换，这个表达式关于 \(a,b\) 对称，取反对称部分后为零，故两种星积给出同一个括号。
\end{proof}

在 Hochschild 余链的语言中，一次结合律就是 \(\delta C_1=0\)，二次结合律则为 \(\delta C_2=C_1\circ C_1\)。命题从这两条方程分别推出乘积法则与 Jacobi 恒等式。这个推论在任意特征下成立，因为括号取的是差 \(C_1(a,b)-C_1(b,a)\)，没有除以 \(2\)。仅当 \(2\) 在 \(k\) 中可逆且选取 \(C_1=\dfrac12\{-,-\}\) 时，才能把一次系数直接写成括号的一半。

\ProseParagraphBreak
星积要求每一阶的结合律，直接求全部 \(C_r\) 通常比寻找一次括号困难。两两交换的导子给出一个可以完整核验的构造：把对两个因子分别求导的算子作形式指数，其三重乘法由同一组三个算子控制。

\begin{proposition}\label{b4:prop:commuting-derivation-star-product}
	设 \(k\) 为特征零域，\(A\) 为交换含幺 \(k\) 代数，\(D_1,\ldots,D_s,E_1,\ldots,E_s\) 为两两交换的 \(k\) 导子。记 \(m:A\otimes_kA\to A\) 为原乘法，并令
	\[
		P=\sum_{i=1}^s D_i\otimes E_i.
	\]
	则下式在 \(A[\![\hbar]\!]\) 上定义形式星积：
	\[
		a\star b=m\exp(\hbar P)(a\otimes b)
		=\sum_{r\ge0}\frac{\hbar^r}{r!}\,mP^r(a\otimes b).
	\]
	它产生的 Poisson 括号为
	\[
		\{a,b\}=\sum_i\bigl(D_i(a)E_i(b)-D_i(b)E_i(a)\bigr).
	\]
\end{proposition}
\begin{proof}
	每个系数只使用有限次导子复合，不需要指数级数在分析意义下收敛。特征零保证所有 \(r!\) 可逆。将公式按 \(\hbar\) 次数延拓，便得到连续双线性乘法。导子满足 \(D_i(1)=E_i(1)=0\)，所以含有单位的高阶项全为零，单位仍为 \(1\)。

	在 \(A^{\otimes3}\) 上记
	\[
		\begin{aligned}
			P_{12}&=\sum_i D_i\otimes E_i\otimes\operatorname{id},\\
			P_{13}&=\sum_i D_i\otimes\operatorname{id}\otimes E_i,\\
			P_{23}&=\sum_i\operatorname{id}\otimes D_i\otimes E_i.
		\end{aligned}
	\]
	这三个算子两两交换；即使它们作用于同一张量因子，该因子上的导子也由假设相互交换。导子的乘积法则给出
	\[
		\begin{aligned}
			P(m\otimes\operatorname{id})
			&=(m\otimes\operatorname{id})(P_{13}+P_{23}),\\
			P(\operatorname{id}\otimes m)
			&=(\operatorname{id}\otimes m)(P_{12}+P_{13}).
		\end{aligned}
	\]
	例如第一式把 \(D_i(ab)\otimes E_i(c)\) 展开为 \(\bigl(D_i(a)\cdot b\bigr)\otimes E_i(c)+\bigl(a\cdot D_i(b)\bigr)\otimes E_i(c)\)。反复应用相同等式，便可将外层的每个 \(P^r\) 移到三重张量上。

	令 \(m_3(a\otimes b\otimes c)=abc\)。两种加括号的乘积因而分别为
	\[
		\begin{aligned}
			(a\star b)\star c
			&=m_3\exp\bigl(\hbar(P_{13}+P_{23})\bigr)
			\exp(\hbar P_{12})(a\otimes b\otimes c),\\
			a\star(b\star c)
			&=m_3\exp\bigl(\hbar(P_{12}+P_{13})\bigr)
			\exp(\hbar P_{23})(a\otimes b\otimes c).
		\end{aligned}
	\]
	交换算子的二项式公式使两式都等于
	\[
		m_3\exp\bigl(\hbar(P_{12}+P_{13}+P_{23})\bigr)(a\otimes b\otimes c).
	\]
	所以乘法结合。这一验证逐个形式系数成立，故也适用于一般形式级数。一次系数为 \(\sum_iD_i(a)E_i(b)\)，取其反对称部分，再用\cref{b4:prop:star-first-commutator-poisson}，便得到所列 Poisson 括号。
\end{proof}

取 \(A=k[q_1,\ldots,q_d,p_1,\ldots,p_d]\)，其中 \(\operatorname{char}k=0\)。所有偏导子相互交换，因此可以取
\[
	P=\frac12\sum_{i=1}^d
	\left(\frac{\partial}{\partial q_i}\otimes\frac{\partial}{\partial p_i}
	-\frac{\partial}{\partial p_i}\otimes\frac{\partial}{\partial q_i}\right).
\]
系数 \(\dfrac12\)、\(-\dfrac12\) 可吸收到各导子中，故它符合刚才的构造。所得括号是\term[biaozhun-xin-kuohao]{标准辛括号}[standard symplectic Poisson bracket]：
\begin{equation}\label{b4:eq:standard-symplectic-poisson}
	\{f,g\}=\sum_{i=1}^d
	\left(\frac{\partial f}{\partial q_i}\frac{\partial g}{\partial p_i}
	-\frac{\partial f}{\partial p_i}\frac{\partial g}{\partial q_i}\right).
\end{equation}
例如 \(\{q_i,p_j\}=\delta_{ij}\)，而两个 \(q\) 或两个 \(p\) 的括号为零。这里 Jacobi 恒等式已由结合星积推出，不必另行排列二阶偏导项。对坐标函数，二阶以上偏导都为零，因而完整的星乘关系就是
\[
	q_i\star p_j=q_ip_j+\frac{\hbar}{2}\delta_{ij},\qquad
	p_j\star q_i=q_ip_j-\frac{\hbar}{2}\delta_{ij},\qquad
	[q_i,p_j]_{\star}=\hbar\delta_{ij}.
\]
一对变量时还可算得
\[
	q^2\star p^2=q^2p^2+2\hbar qp+\frac{\hbar^2}{2},\qquad
	p^2\star q^2=q^2p^2-2\hbar qp+\frac{\hbar^2}{2}.
\]
二次项分别来自 \(\dfrac{\hbar^2}{2!}\cdot\dfrac14\cdot2\cdot2\)，所以两项相减后为 \(4\hbar qp\)，与 \(\{q^2,p^2\}=4qp\) 一致。固定两个多项式时这些求导级数实际终止；对任意形式级数，仍只能按 \(\hbar\) 的次数解释它。此构造使用常系数且相互交换的导子，并不证明任意 Poisson 括号都有星积，也不能在正特征下原样使用带 \(r!\) 分母的公式。


% !TeX root = ../../Book4.tex
% 纲要标题：### 13.5 循环同调与混合复形
% * 状态：初稿已完成，含完整证明、算例及配套习题解答；尚未正式审核。
% #### 13.5.1 Connes 算子与混合复形
% #### 13.5.2 循环复形与 SBI 长正合列
% #### 13.5.3 基域与简单代数的计算
% ---

\section{循环同调与混合复形}
\label{b4:sec:1305}

矩阵迹满足 \(\operatorname{tr}(ab)=\operatorname{tr}(ba)\)。把这个等式反复用于乘积，就得到 \(\operatorname{tr}(a_0\cdots a_n)=\operatorname{tr}(a_n a_0\cdots a_{n-1})\)。Hochschild 链的最后一个边界项也把末因子移到开头再相乘，因此循环换序已经出现在它的公式中。不过，单独旋转一条张量链并不与 Hochschild 边界交换。循环同调以 Hochschild 链复形为基础，再加入一个向上升次的算子，把旋转与边界之间的关系一起记入新的复形。这一构造的混合复形形式见 Weibel\cite[Section 9.8, pp. 344--348]{Weibel1994HomologicalAlgebra}。

\subsection{Connes 算子与混合复形}
\label{b4:subsec:1305:01}

以下仍设 \(k\) 为域，\(A\) 为结合含幺 \(k\) 代数，所有张量积均在 \(k\) 上进行。只取系数双模 \(A\) 本身，记
\[
	C_n(A)=A^{\otimes(n+1)},\qquad C_n(A)=0\quad(n<0).
\]
为使循环换序易于辨认，将纯张量 \(a_0\otimes\cdots\otimes a_n\) 简写为 \((a_0,\ldots,a_n)\)；这只是张量记号，并不把 \(C_n(A)\) 改成笛卡尔积。采用\eqref{b4:eq:hochschild-boundary}的边界 \(b_n:C_n(A)\to C_{n-1}(A)\)，并规定 \(b_0=0\)。这里采用未归一化链：不模去那些在 \(a_1,\ldots,a_n\) 的某个位置插入单位所得到的退化链，后面的公式也保留这些项。

\begin{definition}\label{b4:def:cyclic-chain-operators}
	设 \(k\) 为域，\(A\) 为结合含幺 \(k\) 代数。在 \(C_n(A)=A^{\otimes_k(n+1)}\) 上定义带符号的循环算子、循环和及插入单位的映射
	\[
		\begin{aligned}
			t_n(a_0,\ldots,a_n)&=(-1)^n(a_n,a_0,\ldots,a_{n-1}),\\
			N_n&=1+t_n+\cdots+t_n^n,\\
			s_n(a_0,\ldots,a_n)&=(1,a_0,\ldots,a_n).
		\end{aligned}
	\]
	另令 \(b'_0=0\)，并在 \(n\ge1\) 时定义
	\[
		b'_n(a_0,\ldots,a_n)=\sum_{i=0}^{n-1}(-1)^i(a_0,\ldots,a_i a_{i+1},\ldots,a_n).
	\]
	与 \(b_n\) 相比，\(b'_n\) 少了首尾相乘的最后一项。
\end{definition}
\symidx{\(t_n,N_n\)：Hochschild 链的带符号循环算子与循环和}
\symidx{\(b'_n\)：去掉首尾相乘项的 Hochschild 边界}

旋转 \(n+1\) 次恢复原张量，累计符号为 \((-1)^{n(n+1)}=1\)，故
\[
	t_n^{n+1}=1,\qquad (1-t_n)N_n=N_n(1-t_n)=0.
\]
这里的 \(N_n\) 是求和，未除以 \(n+1\)，所以不要求该整数在 \(k\) 中可逆。最低次有 \(t_0=N_0=1\)，而 \(t_1(a_0,a_1)=-(a_1,a_0)\)。带符号的旋转与普通旋转必须分开。

\begin{lemma}\label{b4:lem:cyclic-boundary-identities}
	设 \(k\) 为域，\(A\) 为结合含幺 \(k\) 代数。上述算子满足
	\begin{equation}\label{b4:eq:cyclic-boundary-identities}
		\begin{aligned}
			b_n(1-t_n)&=(1-t_{n-1})b'_n, & n&\ge1,\\
			b'_nN_n&=N_{n-1}b_n, & n&\ge1,\\
			b'_{n+1}s_n+s_{n-1}b'_n&=1, & n&\ge0,
		\end{aligned}
	\end{equation}
	其中 \(s_{-1}=0\)。
\end{lemma}
\begin{proof}
	记 \(\partial_i\) 为不带符号的第 \(i\) 个面映射：当 \(i<n\) 时相乘第 \(i,i+1\) 个因子，\(\partial_n(a_0,\ldots,a_n)=(a_n a_0,a_1,\ldots,a_{n-1})\)。令 \(\delta_i=(-1)^i \partial_i\)，于是 \(b_n=\sum_{i=0}^n\delta_i\)、\(b'_n=\sum_{i=0}^{n-1}\delta_i\)。先旋转再相乘，可直接读出
	\[
		\delta_0t_n=\delta_n,\qquad \delta_i t_n=t_{n-1}\delta_{i-1}\quad(1\le i\le n).
	\]
	第一式把旋转后的前两个因子 \(a_n,a_0\) 相乘，正好给出原来的末项；第二式把相邻合并的位置往前移一位，两次定义中的符号也恰好一致。因此 \(b_n t_n=\delta_n+t_{n-1}b'_n\)，从 \(b_n=b'_n+\delta_n\) 减去它，得到第一条恒等式。

	第二条需要把所有旋转相加。迭代刚才的面映射公式，有
	\[
		\delta_i t_n^j=
		\begin{cases}
			t_{n-1}^j\delta_{i-j},&i\ge j,\\
			t_{n-1}^{j-1}\delta_{n+1+i-j},&i<j,
		\end{cases}
		\qquad 0\le i\le n-1,\quad 0\le j\le n.
	\]
	第二种情形表示合并位置在旋转中越过了首尾。现在展开 \(b'_nN_n=\sum_{i,j}\delta_i t_n^j\)。固定 \(0\le p\le n-1\)，含 \(t_{n-1}^p\) 的项有两组：第一行取 \(j=p\)，留下 \(\delta_0,\ldots,\delta_{n-1-p}\)；第二行取 \(j=p+1\)，留下 \(\delta_{n-p},\ldots,\delta_n\)。它们合起来正是 \(t_{n-1}^p b_n\)。再对 \(p\) 求和便得 \(b'_nN_n=N_{n-1}b_n\)。

	最后，\(b'_{n+1}s_n(a_0,\ldots,a_n)\) 的第一项将新插入的 \(1\) 与 \(a_0\) 相乘，得到原张量。其余项分别是先在原张量内相邻相乘、再前插 \(1\) 的结果，符号比 \(s_{n-1}b'_n\) 中对应项多一个负号。因此第三式成立；在 \(n=0\) 时，它就是 \(b'_1(1,a_0)=a_0\)。
\end{proof}

相邻两次合并的两种次序由结合律给出相同张量，而符号相反，所以 \((b')^2=0\)，与双模 bar 复形中的计算相同。第三条等式说明这个去掉首尾项的复形可被 \(s\) 收缩。将它与循环和、循环差合成，便能把这份收缩用于原来的 Hochschild 复形。

\begin{definition}\label{b4:def:mixed-complex}
	设 \(k\) 为域，\(C_n\) 为 \(k\) 向量空间，且 \(C_n=0\) 对 \(n<0\) 成立。若线性映射
	\[
		b:C_n\longrightarrow C_{n-1},\qquad B:C_n\longrightarrow C_{n+1}
	\]
	满足 \(b^2=B^2=bB+Bb=0\)，则称 \((C_\bullet,b,B)\) 为一个\term[hunhe-fuxing]{混合复形}[mixed complex]。对结合含幺 \(k\) 代数 \(A\)，定义 \(C_\bullet(A)\) 上的\term[connes-suanzi]{Connes 算子}[Connes operator]为
	\begin{equation}\label{b4:eq:connes-operator}
		B_n=(1-t_{n+1})s_nN_n:C_n(A)\longrightarrow C_{n+1}(A).
	\end{equation}
\end{definition}
\symidx{\(B_n\)：升次的 Connes 算子}

Connes 算子将链的所有循环位置求和，再插入单位并取一次循环差。\footnote{孔涅（Alain Connes，1947-），法国数学家，研究算子代数与非交换几何，并发展循环上同调。}这三个操作的顺序使下列反交换恒等式成立。

\begin{proposition}\label{b4:prop:hochschild-mixed-complex}
	设 \(k\) 为任意特征的域，\(A\) 为结合含幺 \(k\) 代数。则 Hochschild 边界 \(b\) 与\eqref{b4:eq:connes-operator}的 \(B\) 使 \(C_\bullet(A)\) 成为混合复形。
\end{proposition}
\begin{proof}
	\Cref{b4:prop:algebra-bar-resolution}及其后的张量计算已证明 \(b^2=0\)。由 \(N_{n+1}(1-t_{n+1})=0\)，有
	\[
		B_{n+1}B_n=(1-t_{n+2})s_{n+1}N_{n+1}(1-t_{n+1})s_nN_n=0.
	\]
	再依次使用\eqref{b4:eq:cyclic-boundary-identities}，得到
	\[
		\begin{aligned}
			b_{n+1}B_n+B_{n-1}b_n
			&=(1-t_n)\bigl(b'_{n+1}s_n+s_{n-1}b'_n\bigr)N_n\\
			&=(1-t_n)N_n=0.
		\end{aligned}
	\]
	在 \(n=0\) 时，第二个复合为零，而同一计算以 \(1-t_0=0\) 结束。因此所有次数上都有所需的反交换关系。
\end{proof}

例如前两个 Connes 算子为
\[
	\begin{aligned}
		B_0(a)&=(1,a)+(a,1),\\
		B_1(a,c)&=(1,a,c)-(1,c,a)-(c,1,a)+(a,1,c).
	\end{aligned}
\]
第一式的第二项前是正号，因为 \(t_1\) 本身带有负号。又有 \(b_2(a,1,1)=(a,1)\)，故这一项在 Hochschild 同调中为零。若 \(A\) 交换，\cref{b4:prop:hh1-differentials}把 \([(a,c)]\) 送到 \(a\cdot\mathrm{d}c\)，因而 \([B_0(a)]\) 对应 \(\mathrm{d}a\)。这说明升次算子 \(B\) 在最低次恢复了熟悉的微分。

\subsection{循环复形与 SBI 长正合列}
\label{b4:subsec:1305:02}

要把降次的 \(b\) 与升次的 \(B\) 放进同一个链复形，需要把 Hochschild 链重复放置，并让每一份相隔两个总次数。这样 \(B:C_{n-2}\to C_{n-1}\) 就能与 \(b:C_n\to C_{n-1}\) 一起成为总边界的两个分量。

\begin{definition}\label{b4:def:cyclic-complex}
	设 \(k\) 为域，\((C_\bullet,b,B)\) 为非负次数的混合复形。定义其\term[xunhuan-fuxing]{循环复形}[cyclic complex]为
	\[
		\operatorname{CC}_n(C)=\bigoplus_{j=0}^{\lfloor n/2\rfloor}C_{n-2j}\quad(n\ge0),\qquad \operatorname{CC}_n(C)=0\quad(n<0),
	\]
	其微分规定为
	\begin{equation}\label{b4:eq:cyclic-total-differential}
		\mathrm{d}(x_0,x_1,\ldots)=(b x_0+B x_1,\ b x_1+B x_2,\ldots),
	\end{equation}
	其中 \(x_j\in C_{n-2j}\)，超出指标范围的分量一律为零。对结合含幺 \(k\) 代数 \(A\)，用 \((C_\bullet(A),b,B)\) 得到 \(\operatorname{CC}_\bullet(A)\)，并定义它的\term[xunhuan-tongdiao]{循环同调}[cyclic homology]为
	\[
		\operatorname{HC}_n(A)=H_n(\operatorname{CC}_\bullet(A)).
	\]
\end{definition}
\symidx{\(\operatorname{CC}_\bullet(A),\operatorname{HC}_n(A)\)：循环复形与循环同调}

每个总次数只含有限项。对第 \(j\) 个分量，两次微分给出
\[
	b^2x_j+(bB+Bb)x_{j+1}+B^2x_{j+2}=0,
\]
所以这确为链复形。其开头是
\[
	\cdots\longrightarrow C_3\oplus C_1
	\xrightarrow{\ \left(\begin{smallmatrix}b&B\\0&b\end{smallmatrix}\right)\ }
	C_2\oplus C_0\xrightarrow{\ (b\ \ B)\ }C_1\xrightarrow{\ b\ }C_0\longrightarrow0.
\]
例如二次循环为 \((x_0,x_1)\in C_2\oplus C_0\)，条件是 \(b x_0=-B x_1\)。因此 \(x_0\) 未必单独是 Hochschild 循环；低两次的分量 \(x_1\) 可以通过 \(B\) 补偿它的边界。这正是把两种算子同时保留的意义。代数同态逐因子作用于张量，与 \(b,t,N,s\) 都交换，因而诱导循环复形的链映射以及循环同调上的映射。

\begin{theorem}[SBI 长正合列]\label{b4:thm:cyclic-sbi}
	设 \(k\) 为域，\(A\) 为结合含幺 \(k\) 代数。则有关于 \(A\) 自然的长正合列
	\[
		\begin{aligned}
			\cdots&\longrightarrow\operatorname{HC}_{n+1}(A)
			\xrightarrow{S}\operatorname{HC}_{n-1}(A)
			\xrightarrow{B}\operatorname{HH}_n(A)\\
			&\xrightarrow{I}\operatorname{HC}_n(A)
			\xrightarrow{S}\operatorname{HC}_{n-2}(A)
			\xrightarrow{B}\operatorname{HH}_{n-1}(A)\longrightarrow\cdots.
		\end{aligned}
	\]
	负次同调群取零。这里 \(I\) 由包含首分量诱导，\(S\) 由删去首分量诱导。若 \((y_0,y_1,\ldots)\) 代表 \(\operatorname{HC}_{n-1}(A)\) 中的类，则连接同态将它送到 \([B y_0]\in\operatorname{HH}_n(A)\)。
\end{theorem}
\symidx{\(S,B,I\)：循环同调与 Hochschild 同调间的长正合列映射}
\begin{proof}
	定义链映射
	\[
		\begin{aligned}
			I_n:C_n(A)&\longrightarrow\operatorname{CC}_n(A),&
			x&\longmapsto(x,0,\ldots),\\
			S_n:\operatorname{CC}_n(A)&\longrightarrow\operatorname{CC}_{n-2}(A),&
			(x_0,x_1,x_2,\ldots)&\longmapsto(x_1,x_2,\ldots).
		\end{aligned}
	\]
	\eqref{b4:eq:cyclic-total-differential}给出 \(\mathrm{d}I=Ib\) 与 \(S\mathrm{d}=\mathrm{d}S\)。这里把 \(\operatorname{CC}_{\bullet-2}(A)\) 看作第 \(n\) 项为 \(\operatorname{CC}_{n-2}(A)\)、微分仍为 \(\mathrm d\) 的链复形；平移两次不改变微分符号。\(I\) 单射，\(S\) 满射，且 \(\ker S\) 恰为首分量，所以有短正合列
	\[
		0\longrightarrow(C_\bullet(A),b)
		\xrightarrow{I}\operatorname{CC}_\bullet(A)
		\xrightarrow{S}\operatorname{CC}_{\bullet-2}(A)\longrightarrow0.
	\]
	将链指标改写为上链指标 \(C^q=C_{-q}\)，便可应用\cref{b4:thm:cohomology-long-exact}；再换回链指标，得到所述长正合列。

	具体地，\(y=(y_0,y_1,\ldots)\) 是 \(\operatorname{CC}_{n-1}(A)\) 的循环，在短正合列的右端它位于第 \(n+1\) 次。选取提升 \(\widetilde y=(0,y_0,y_1,\ldots)\)，则
	\[
		\mathrm d\widetilde y=(B y_0,\ b y_0+B y_1,\ldots)=(B y_0,0,\ldots)=I(B y_0).
	\]
	这就是连接同态的公式。由 \(bB=-Bb\) 及 \(b y_0=-B y_1\)，有 \(bB y_0=B^2y_1=0\)，所以它确为 Hochschild 循环。改变循环代表时，若增加 \(\mathrm dz\)，首分量增加 \(b z_0+B z_1\)，其 \(B\) 像增加 \(-bB z_0\)，是 Hochschild 边界。代数同态与提升后取边界的过程相容，故整条长正合列自然。
\end{proof}

\(S\) 降低两次，但一般并不是同构。SBI 列的低次部分为
\[
	\operatorname{HH}_2(A)\longrightarrow\operatorname{HC}_2(A)
	\xrightarrow{S}\operatorname{HC}_0(A)
	\xrightarrow{B}\operatorname{HH}_1(A)
	\xrightarrow{I}\operatorname{HC}_1(A)\longrightarrow0.
\]
零次没有新增边界，故 \(\operatorname{HC}_0(A)=\operatorname{HH}_0(A)=A/[A,A]\)，其中 \([A,A]\) 是交换子张成的 \(k\) 子空间。长正合列说明一次循环同调进一步模去了由 \(B\) 产生的类。它也解释了名称中的三个字母：\(S\) 记录相隔两次的重复，\(I\) 把 Hochschild 类送入循环同调，\(B\) 测量一个循环类能否向高两次提升。

\begin{corollary}\label{b4:cor:cyclic-first-differentials}
	设 \(k\) 为域，\(A\) 为交换含幺 \(k\) 代数。则
	\[
		\operatorname{HC}_0(A)=A,\qquad
		\operatorname{HC}_1(A)\cong\Omega_{A/k}/\mathrm dA
	\]
	是自然的 \(k\) 向量空间同构，其中 \(\mathrm dA\) 指导子 \(\mathrm d:A\to\Omega_{A/k}\) 的像。
\end{corollary}
\begin{proof}
	交换性给出 \([A,A]=0\)。由\cref{b4:prop:hh1-differentials}，\(\operatorname{HH}_1(A)\cong\Omega_{A/k}\)。在这个同构下，SBI 列中的 \(B:A\to\operatorname{HH}_1(A)\) 把 \(a\) 送到 \(\mathrm d a\)，因为 \(B_0(a)=(1,a)+(a,1)\) 对应 \(\mathrm d a+a\cdot\mathrm d1=\mathrm d a\)。于是取低次列的余核便得结论。
\end{proof}

这里的 \(\mathrm dA\) 一般只是 \(k\) 子空间，未必是 \(A\) 子模，故不能把该商自动当作 \(A\) 模。虽然混合复形的构造对任意特征都成立，求导的像却会随特征改变。例如在特征 \(p>0\) 时，\(\mathrm d(x^p)=0\)，这会影响 \(S\) 的核以及一次循环同调的计算。

\subsection{基域与简单代数的计算}
\label{b4:subsec:1305:03}

SBI 列将一类循环同调计算变成已有 Hochschild 计算的直接后果。若正次 Hochschild 同调全部消失，连接同态两侧便没有剩余障碍，\(S\) 就在非负次数间给出二周期同构。重要的是先验证这一消失条件，而不是仅由代数为简单环就作推断。

\begin{proposition}\label{b4:prop:cyclic-hh-vanishing}
	设 \(k\) 为域，\(A\) 为结合含幺 \(k\) 代数，并满足 \(\operatorname{HH}_n(A)=0\) 对所有 \(n>0\) 成立。则对每个 \(q\ge0\)，有自然同构
	\[
		\operatorname{HC}_{2q}(A)\cong A/[A,A],\qquad
		\operatorname{HC}_{2q+1}(A)=0.
	\]
	偶数次的同构由 \(S^q:\operatorname{HC}_{2q}(A)\to\operatorname{HC}_0(A)\) 给出。
\end{proposition}
\begin{proof}
	\(\operatorname{HC}_0(A)=A/[A,A]\)，且 \(\operatorname{HH}_1(A)\to\operatorname{HC}_1(A)\to0\) 正合，故 \(\operatorname{HC}_1(A)=0\)。对 \(n\ge2\)，SBI 列包含
	\[
		0=\operatorname{HH}_n(A)\longrightarrow\operatorname{HC}_n(A)
		\xrightarrow{S}\operatorname{HC}_{n-2}(A)
		\longrightarrow\operatorname{HH}_{n-1}(A)=0.
	\]
	因而 \(S\) 是同构，逐次降至零次或一次即得结果。
\end{proof}

对基域 \(A=k\)，包络代数也是 \(k\)，所以 \(\operatorname{HH}_n(k)=\operatorname{Tor}^k_n(k,k)=0\) 对 \(n>0\) 成立，\(\operatorname{HH}_0(k)=k\)。于是
\[
	\operatorname{HC}_{2q}(k)=k,\qquad\operatorname{HC}_{2q+1}(k)=0.
\]
这在任意特征下都成立。未归一化模型的 \(C_n(k)\) 在每个次数上都是 \(k\)，却不表示其所有正次同调非零：在这些识别下，\(b_n\) 在正偶数次为恒等，在奇数次为零，正次同调确实全部消失。

\begin{definition}\label{b4:def:separable-algebra-bimodule}
	设 \(k\) 为域，\(A\) 为有限维结合含幺 \(k\) 代数。若乘法满射
	\[
		\mu:A\otimes_k A\longrightarrow A,\qquad a\otimes c\longmapsto ac
	\]
	作为 \(A\) 双模同态有一个双模截面，则称 \(A\) 为\term[kefen-daishu]{可分代数}[separable algebra]。这里 \(A\otimes_k A\) 的左右作用分别在最外侧因子上。
\end{definition}

可分性在这里落实为一份可以计算的双模截面。截面 \(\rho\) 完全由 \(e=\rho(1)=\sum_i u_i\otimes v_i\) 决定，条件等价于
\begin{equation}\label{b4:eq:separability-element}
	\sum_i u_i v_i=1,\qquad
	\sum_i a u_i\otimes v_i=\sum_i u_i\otimes v_i a\quad(a\in A).
\end{equation}
事实上，截面给出这两个等式；反过来，若 \(e\) 满足它们，则 \(\rho(a)=\sum_i a u_i\otimes v_i\) 同时保持左右作用，并满足 \(\mu\rho=1_A\)。有限维可分性的双模表述见 Weibel\cite[Theorem 9.2.11, pp. 310--311]{Weibel1994HomologicalAlgebra}。下面直接从截面得到同调的消失。

\begin{proposition}\label{b4:prop:separable-cyclic-homology}
	设 \(k\) 为域，\(A\) 为有限维可分 \(k\) 代数。则 \(\operatorname{HH}_n(A)=0\) 对 \(n>0\) 成立，且
	\[
		\operatorname{HC}_{2q}(A)\cong A/[A,A],\qquad
		\operatorname{HC}_{2q+1}(A)=0\qquad(q\ge0).
	\]
\end{proposition}
\begin{proof}
	带有外侧作用的 \(A\otimes_k A\) 是秩一自由左 \(A^e\) 模：\(a\otimes c^{\mathrm{op}}\mapsto a\otimes c\) 给出同构。乘法映射的双模截面使 \(A\) 成为这个自由模的直和因子，因此 \(A\) 是投射左 \(A^e\) 模。计算 \(\operatorname{Tor}^{A^e}_n(A^{\mathrm r},A)\) 时可以用只含第零次 \(A\) 的投射消解，故正次 Hochschild 同调为零。再应用\cref{b4:prop:cyclic-hh-vanishing}即可。
\end{proof}

现在令 \(A=M_r(k)\)，其中 \(r\ge1\)。用矩阵单位写
\[
	e=\sum_{i=1}^r E_{i1}\otimes E_{1i}.
\]
一方面 \(\mu(e)=\sum_iE_{ii}=1_A\)；另一方面，对每个矩阵单位 \(E_{uv}\)，\eqref{b4:eq:separability-element}两边都等于 \(E_{u1}\otimes E_{1v}\)。由线性，这个等式对每个矩阵成立，故 \(M_r(k)\) 可分。这里没有除以 \(r\)，所以即使 \(\operatorname{char}k\) 整除 \(r\)，论证仍成立。

\ProseParagraphBreak
还需识别 \(A/[A,A]\)。普通矩阵迹消去交换子，且每个非对角矩阵单位和对角差都满足
\[
	E_{ij}=[E_{ii},E_{ij}]\quad(i\ne j),\qquad
	E_{ii}-E_{jj}=[E_{ij},E_{ji}].
\]
这些矩阵张成所有迹为零的矩阵：非对角部分已经包含，而迹为零的对角矩阵可写成 \(\sum_{i<r}a_{ii}(E_{ii}-E_{rr})\)。再由 \(\operatorname{tr}(E_{11})=1\)，得到 \(A/[A,A]\xrightarrow{\sim}k\)。于是
\[
	\operatorname{HC}_{2q}(M_r(k))\cong k,\qquad
	\operatorname{HC}_{2q+1}(M_r(k))=0.
\]
零次同构由普通迹给出，所有偶数次同构则再复合 \(S^q\)。当 \(r\cdot1_k=0\) 时，单位矩阵的交换子商类为零，生成这个一维商空间的是 \([E_{11}]\)，不能用除以 \(r\) 的归一化迹代替。这也与\cref{b4:thm:hh-matrix-morita}所得的 Hochschild 消失相符。

\ProseParagraphBreak
有限维的限制本身仍不足以把简单性变成可分性。例如令 \(k=\mathbb F_p(t)\)、\(L=k(u)\)，其中 \(u^p=t\)。\(U^p-t\) 对 \(\mathbb F_p[t]\) 中的素元 \(t\) 满足 Eisenstein 判据，故在 \(k[U]\) 中不可约，\([L:k]=p\)。\(L\) 是域，因而是简单 \(k\) 代数，但这个扩张纯不可分。由 \(L=k[U]/(U^p-t)\) 计算微分模时，关系的微分为 \(p u^{p-1}\cdot\mathrm d u=0\)，没有消去 \(\mathrm d u\)，所以 \(\Omega_{L/k}=L\,\mathrm d u\ne0\)。也可直接检验非零性：\(k[U]\) 上的形式导数保持理想 \((U^p-t)\)，因而下降为 \(L\) 上的 \(k\) 导子，并将 \(u\) 送到 \(1\)。若 \(\mathrm d u=0\)，微分模的泛性质便迫使这个导子在 \(u\) 上也取零值，矛盾。\Cref{b4:prop:hh1-differentials}给出 \(\operatorname{HH}_1(L)\ne0\)，因此前面的消失命题不适用。循环同调的计算始终依赖所选基域；对同一个域 \(L\)，相对于 \(L\) 自身计算和相对于 \(k\) 计算，所得复形不同。


\begin{chaptersummary}
	Hochschild 理论把乘法问题放入双模范畴：bar 消解提供明确计算，中心、导子和微分解释低次数，二次上同调记录一阶乘法改动在坐标变换下的等价类。等价要求模参数还原为恒等映射，这是分类结论的一部分。

	\ProseParagraphBreak
	把结合律展开到更高阶，所得障碍是具体的三余圈。三次上同调全消失足以保证形式延拓，却不是延拓存在的必要条件；非零二次类也不保证可以延拓。矩阵代数的不变性依靠双模范畴等价，多项式代数的计算依靠 Koszul 消解，光滑比较的概要则说明局部正则对角怎样给出同一类消解。三种方法分别适合不同的代数，不能仅凭公式相似互相替换。

	\ProseParagraphBreak
	当原代数交换时，形式星积的交换子中参数的一次项系数给出 Poisson 括号。一阶结合律只保证双导子性质，Jacobi 恒等式还需要下一阶结合律；常系数偏导构造的星积展示了怎样把这种低阶数据延拓为所有阶的乘法。循环同调则保留 Hochschild 链上的另一个算子 \(B\)：关系 \(bB+Bb=0\) 与 \(B^2=0\) 使多个链次数能组成一个循环复形。移去它的第一列所得到的短正合列产生 SBI 长正合列；当正次 Hochschild 同调消失时，这条列将循环同调归结为零次交换子商及其偶次重复。
\end{chaptersummary}

\BookExercises

\begin{optionalexercise}\label{b4:ex:12-matrix-envelope}
	设 \(k\) 为域，\(r\ge1\)，\(A=M_r(k)\)，\(A^e=A\otimes_kA^{\mathrm{op}}\)。证明 \(A^e\to\operatorname{End}_k(A)\)、\(a\otimes b^{\mathrm{op}}\mapsto(x\mapsto axb)\) 是代数同构。写出矩阵单位作用，并解释为何不需要假设特征不整除 \(r\)。
\end{optionalexercise}
\begin{solution}
	乘法相容性正是包络代数的乘法定义。\(E_{ij}\otimes E_{uv}^{\mathrm{op}}\) 把 \(E_{pq}\) 送到 \(\delta_{jp}\delta_{qu}E_{iv}\)。这个算子只把基向量 \(E_{ju}\) 送到 \(E_{iv}\)，把其他基向量送到零；相对于以 \(E_{pq}\) 为基的坐标，它就是行指标为 \((i,v)\)、列指标为 \((j,u)\) 的矩阵单位。随着四个指标变化，所有输入、输出基向量对恰好各出现一次，因而这些算子构成 \(\operatorname{End}_k(A)\) 的基。原映射于是双射。计算只用系数 \(0,1\)，没有除以 \(r\)。
\end{solution}

\begin{optionalexercise}\label{b4:ex:12-upper-triangular}
	设 \(k\) 为域，\(r\ge2\)，\(T_r\subseteq M_r(k)\) 为上三角矩阵代数。计算 \(\operatorname{HH}^0(T_r)\)、\(\operatorname{HH}_0(T_r)\)，以及包含映射在零次同调上诱导的映射与其核。若 \(\operatorname{char}k\) 整除 \(r\)，比较 \(I_r\) 在两个零次同调空间中的类。
\end{optionalexercise}
\begin{solution}
	与全部 \(E_{ii}\) 交换迫使中心元素为对角矩阵，再与 \(E_{ij}\)、\(i<j\) 交换迫使各对角元相等，故 \(\operatorname{HH}^0(T_r)=kI_r\)。交换子的对角项均为零，且 \([E_{ii},E_{ij}]=E_{ij}\) 对 \(i<j\) 成立，所以 \([T_r,T_r]\) 是严格上三角矩阵空间。因而 \(\operatorname{HH}_0(T_r)\cong k^r\)，由对角元给出。包含映射将交换子送到交换子，故在商上诱导映射；按矩阵代数的迹识别 \(\operatorname{HH}_0(M_r(k))\cong k\)，该映射为
	\[
		(c_1,\ldots,c_r)\longmapsto\sum_{i=1}^r c_i.
	\]
	其核为 \(\{(c_i)\in k^r\mid\sum_i c_i=0\}\)，维数为 \(r-1\)。源中的 \([I_r]\) 对应 \((1,\ldots,1)\)，总是非零；若特征整除 \(r\)，目标中的 \([I_r]\) 却因迹为零而消失。不能把包含映射误当作零次同调的单射。
\end{solution}

\begin{optionalexercise}\label{b4:ex:12-dual-derivation-characteristic}
	设 \(k\) 为域，\(A=k[x]/(x^2)\)。分别在特征二和非二时计算 \(\operatorname{HH}^1(A)\)、\(\operatorname{HH}^2(A)\)，并标出特征条件进入的位置。
\end{optionalexercise}
\begin{solution}
	导子由 \(D(x)=a+bx\) 决定，关系要求 \(0=D(x^2)=2ax\)。非二时 \(a=0\)，故一次上同调为 \(kx\)；特征二时 \(a,b\) 任意，故为 \(A\)。内导子因交换性为零。归一化二次上链由 \(\phi(x,x)\in A\) 决定且全为余圈，对归一化一次上链 \(h\)，余边公式给出 \(\delta h(x,x)=x\cdot h(x)-h(x^2)+h(x)\cdot x=2x\cdot h(x)\)，因为 \(x^2=0\)、\(h(0)=0\)，而 \(A\) 交换。若 \(h(x)=a+bx\)，这个值为 \(2ax\)：非二时余边恰为 \(kx\)，特征二时则全为零。所以二次上同调为 \(A/(2xA)\)：非二时为 \(k\)，特征二时为 \(A\)。两处差别都来自同一个系数 \(2\)。
\end{solution}

\begin{optionalexercise}\label{b4:ex:12-free-algebra-derivation}
	设 \(k\) 为域，\(A=k\langle x,y\rangle\) 为自由结合 \(k\) 代数。给定 \(u,v\in A\)，证明存在唯一 \(k\) 导子 \(D\) 使 \(D(x)=u,D(y)=v\)。证明 \(D(x)=1,D(y)=0\) 不是内导子。
\end{optionalexercise}
\begin{solution}
	对单词逐个位置替换：\(D(z_1\cdots z_n)=\sum_i z_1\cdots z_{i-1}D(z_i)z_{i+1}\cdots z_n\)，并令 \(D(1)=0\)，线性延拓。拼接两个单词时，把求和分为落在左、右单词中的两部分便得 Leibniz 公式；反过来 Leibniz 公式也强制上述定义。若指定的导子为内导子，则 \(1=xm-mx\) 对某个 \(m\) 成立。送入交换化商 \(k[x,y]\)，右边为零，左边为 \(1\)，矛盾。
\end{solution}

\begin{optionalexercise}\label{b4:ex:12-square-zero-extension}
	设 \(k\) 为域，\(A=k[x]/(x^2)\)，\(N=k\) 带有由 \(x\mapsto0\) 得到的 \(A\) 双模作用。把 \(E=k[z]/(z^3)\) 映到 \(A\)，\(z\mapsto x\)，其核用 \(z^2\leftrightarrow1\in N\) 识别。证明核平方零，并证明这个扩张不存在保持单位的代数截面。
\end{optionalexercise}
\begin{solution}
	核为 \(kz^2\)，其平方为零，左右作用都通过常数项。向量空间截面 \(1\mapsto1,x\mapsto z\) 的乘法误差只在 \((x,x)\) 上非零，值为 \(1\in N\)。任何代数截面必须把 \(x\) 送到 \(z+c z^2\)，但其平方为 \(z^2\ne0\)，故不存在。也可用余边检验：对任意 \(h:A\to N\)，\(\delta h(x,x)=x\cdot h(x)-h(0)+h(x)\cdot x=0\)，不能消去所给误差。
\end{solution}

\begin{optionalexercise}\label{b4:ex:12-coordinate-change}
	设 \(k\) 为特征不为二的域，\(A=k[x]/(x^2)\)，\(v\in k\)。比较 \(A\otimes_k k[t]/(t^2)\) 上的一阶乘法 \(x*_\phi x=t(1+vx)\) 与 \(x*_\psi x=t\)，两者都保持原单位。写出还原为恒等的显式等价。
\end{optionalexercise}
\begin{solution}
	取 \(h(1)=0,h(x)=\dfrac v2\)，则 \(\delta h(x,x)=vx=\phi(x,x)-\psi(x,x)\)。由分类证明，\(T=1+th\) 是从 \(*_\phi\) 到 \(*_\psi\) 的代数同构。直接检验：\(T(x*_\phi x)=t+tvx\)，而在目标乘法下 \(\bigl(x+t\dfrac v2\bigr)*\bigl(x+t\dfrac v2\bigr)=t+tvx\)，单位也保持。
\end{solution}

\begin{optionalexercise}\label{b4:ex:12-noncommutative-deformation}
	设 \(k\) 为域，\(A=k[x,y]\)，令 \(\phi(a,b)=\dfrac{\partial a}{\partial x}\dfrac{\partial b}{\partial y}\)。证明它定义保持单位的一阶形变，其上同调类非零，且所得形变不交换。
\end{optionalexercise}
\begin{solution}
	两个偏导都是一次余圈，所以\cref{b4:prop:hochschild-cup-differential}说明其杯积 \(\phi\) 是\(2\)-余圈；偏导零化单位，故归一化。\(\phi(x,y)=1\)，\(\phi(y,x)=0\)，因此 \(x*y-y*x=t\)。交换代数上每个\(2\)-余边 \(\delta h(a,b)=a\cdot h(b)-h(ab)+h(a)\cdot b\) 都关于 \(a,b\) 对称，\(\phi\) 却不对称，故不是余边。这个结论在特征二也成立。
\end{solution}

\begin{optionalexercise}\label{b4:ex:12-derivation-bracket}
	设 \(k\) 为域，\(A=k[x]\)，\(f,g\in k[x]\)。取 \(D=f(x)\cdot\dfrac{\mathrm{d}}{\mathrm{d}x}\)、\(E=g(x)\cdot\dfrac{\mathrm{d}}{\mathrm{d}x}\)。计算它们在 \(\operatorname{HH}^1(A)\) 中的括号，并给出一次非零的具体计算。
\end{optionalexercise}
\begin{solution}
	对任意 \(a\)，复合差为 \(f(g'a'+ga'')-g(f'a'+fa'')=(fg'-gf')a'\)。故 \([D,E]=(fg'-gf')\cdot\dfrac{\mathrm{d}}{\mathrm{d}x}\)。取 \(f=1,g=x\)，得到 \(\bigl[\dfrac{\mathrm{d}}{\mathrm{d}x},x\cdot\dfrac{\mathrm{d}}{\mathrm{d}x}\bigr]=\dfrac{\mathrm{d}}{\mathrm{d}x}\)，在任何特征都非零。内导子为零，所以这也正是上同调中的非零括号。
\end{solution}

\begin{optionalexercise}\label{b4:ex:12-cup-derivation-square}
	设 \(k\) 为特征不为二的域，\(A\) 为结合含幺 \(k\) 代数，\(D:A\to A\) 为 \(k\) 导子。证明 \(D\smile D\) 是\(2\)-余边，并写出原像。
\end{optionalexercise}
\begin{solution}
	两次使用 Leibniz 公式得 \(D^2(ab)=D^2(a)b+2D(a)D(b)+aD^2(b)\)。因此 \(\delta(D^2)(a,b)=-2D(a)D(b)\)，所以 \(D\smile D=\delta\bigl(-\dfrac12D^2\bigr)\)。这里不能删掉特征条件后保留同一个原像公式：特征二时，上式只说明 \(\delta(D^2)=0\)，即复合 \(D^2\) 本身也是导子，并未说明 \(D\smile D\) 为余边。
\end{solution}

\begin{optionalexercise}\label{b4:ex:12-formal-rigidity}
	设 \(k\) 为域，\(A\) 为非零结合含幺 \(k\) 代数且 \(\operatorname{HH}^2(A)=0\)。证明每个保持单位的形式形变，都能通过模 \(t\) 为恒等的连续 \(k[\![t]\!]\) 线性坐标变换化为原乘法；连续性相对于 \(t\)-进拓扑。
\end{optionalexercise}
\begin{solution}
	若已经消去次数小于 \(r\) 的改动，乘法形如 \(m+t^r\phi_r\pmod {t^{r+1}}\)。结合律在 \(r\) 次给出 \(\delta\phi_r=0\)，单位使其归一化。二次上同调为零，故存在 \(h_r\) 使 \(\phi_r=\delta h_r\)；归一化迫使 \(h_r(1)=0\)。用 \(T_r=1+t^rh_r\) 从当前乘法搬运到新的乘法，\(r\) 次项即变为零，且不影响更低次数。令 \(T^{(N)}=T_N\circ\cdots\circ T_1\)。因 \(T_{N+1}\) 模 \(t^{N+1}\) 为恒等，\(T^{(N+1)}\) 与 \(T^{(N)}\) 在这一截断中相同；其逆也形成相容系统。因此各有限截断上的代数同构及其逆可以分别取逆极限，得到连续的形式自同构。最终每个正次数项均在某个有限步骤中被消去。这个结论关于形式乘法，不断言任何解析收敛性。
\end{solution}

\begin{optionalexercise}\label{b4:ex:12-automorphism-orbits}
	设 \(k\) 为特征不为二的域，\(A=k[x]/(x^2)\)，\(u\in k\)、\(\lambda\in k^\times\)。在 \(\operatorname{HH}^2(A)\cong k\) 的坐标 \(\phi(x,x)=u\) 下，计算代数自同构 \(g(x)=\lambda x\) 对形变类的作用。若允许等价模 \(t\) 为任意自同构，分类怎样变化？
\end{optionalexercise}
\begin{solution}
	搬运后的上链为 \(\phi'(a,b)=g\phi(g^{-1}a,g^{-1}b)\)。因 \(g^{-1}x=\lambda^{-1}x\)，两个输入各贡献一个 \(\lambda^{-1}\)，而 \(g\) 固定标量 \(u\)，所以
	\[
		\phi'(x,x)=g\bigl(\lambda^{-2}\phi(x,x)\bigr) =\lambda^{-2}u.
	\]
	因此原来的参数集合 \(k\) 要再按非零平方倍数取轨道：零单独成轨道，非零参数按 \(k^\times/(k^\times)^2\) 分类。这个附加商来自放宽还原识别，并非 \(\operatorname{HH}^2\) 定义中已经做过的商。
\end{solution}

\begin{optionalexercise}\label{b4:ex:12-matrix-polynomial}
	设 \(k\) 为域，\(B=M_2(k[x,y])\)。计算所有次数的 \(\operatorname{HH}_n(B)\) 与 \(\operatorname{HH}^n(B)\) 作为 \(k[x,y]\) 模的秩。特征二时零次同调能否由单位矩阵的类生成？
\end{optionalexercise}
\begin{solution}
	矩阵不变性与多项式 Koszul 计算给出次数 \(0,1,2\) 的秩依次为 \(1,2,1\)，更高次数为零。特征二时 \(I_2\) 的迹为零，所以其零次同调类为零，不能生成；\(E_{11}\) 的类对应常数 \(1\)，才给出自由生成元。不能以单位的类是否非零判断整个零次同调是否消失。
\end{solution}

\begin{optionalexercise}\label{b4:ex:12-dual-periodic-hh}
	设 \(k\) 为域，\(A=k[x]/(x^2)\)。在 \(S=A^e=k[x,y]/(x^2,y^2)\) 上，证明由乘 \(x-y\)、乘 \(x+y\) 交替组成、最右一条微分为乘 \(x-y\) 的自由复形消解 \(A\)。由此计算全部正次 Hochschild 同调与上同调。
\end{optionalexercise}
\begin{solution}
	\(S\) 的基为 \(1,x,y,xy\)。令 \(w=a+bx+cy+dxy\)，则
	\[
		w(x-y)=a(x-y)+(c-b)xy,\qquad w(x+y)=a(x+y)+(b+c)xy.
	\]
	两式为零分别要求 \(a=0,c=b\) 与 \(a=0,c=-b\)。因而
	\[
		\operatorname{ann}_S(x-y)=k(x+y)\oplus kxy=(x+y),\qquad \operatorname{ann}_S(x+y)=k(x-y)\oplus kxy=(x-y).
	\]
	这些等式在特征二仍成立。最右商为 \(S/(x-y)=A\)，故复形是自由消解。与 \(A\) 张量后，微分从一次到零次起交替为 \(0,2x,0,2x,\ldots\)。于是对 \(r\ge0\) 以及 \(s\ge1\)，
	\[
		\operatorname{HH}_{2r+1}(A)=A/(2xA),\qquad \operatorname{HH}_{2s}(A)=\operatorname{ann}_A(2x).
	\]
	取 Hom 后上链微分从零次起交替为 \(0,2x,0,2x,\ldots\)，所以
	\[
		\operatorname{HH}^{2r+1}(A)=\operatorname{ann}_A(2x),\qquad \operatorname{HH}^{2s}(A)=A/(2xA).
	\]
	非二特征时两种模分别为 \(kx\) 与 \(k\)；特征二时全部为 \(A\)。这里同调与上同调的奇偶角色不能互换。
\end{solution}

\begin{optionalexercise}\label{b4:ex:12-nonzero-hh3-extension}
	设 \(k\) 为域，\(A=k[x]/(x^2)\)，固定 \(u,v\in k\)。在 \(A\otimes_k k[t]/(t^3)\) 上，保持单位、模 \(t^2\) 为 \(x*x=t(u+vx)\) 的乘法可写成
	\[
		x*x=t(u+vx)+t^2(a+bx),\qquad a,b\in k.
	\]
	证明每对 \((a,b)\) 都给出结合乘法，并分类这些乘法在保持单位、模 \(t^2\) 为恒等的 \(k[t]/(t^3)\) 线性代数同构下的等价类。分别说明特征二与非二的结果。
\end{optionalexercise}
\begin{solution}
	对任意 \(a,b\)，取
	\[
		k[t]/(t^3)[X]/\bigl(X^2-t(u+vX)-t^2(a+bX)\bigr).
	\]
	首一除法使它在 \(k[t]/(t^3)\) 上以 \(1,X\) 为自由基，给出所要求的结合乘法。允许的坐标变换为 \(T=1+t^2h\)，其中 \(h(1)=0\)。设 \(h(x)=c+dx\)，比较源参数 \((a,b)\) 与目标参数 \((a',b')\) 的乘法，二阶系数满足
	\[
		(a-a')+(b-b')x=(\delta h)(x,x)=2cx.
	\]
	故等价当且仅当 \(a=a'\)、\(b-b'\in2k\)。特征不为二时，每个类由 \(a\) 唯一标记，\(b\) 可消去；从 \((a,b)\) 到 \((a,0)\) 的等价可取 \(T(x)=x+t^2\dfrac b2\)。特征二时，等价要求两对参数完全相等。这一分类固定了一阶乘法和到一阶的坐标识别，比仅证明某个形式延拓存在更细。
\end{solution}

\begin{optionalexercise}\label{b4:ex:12-bar-twisted-coefficients}
	设 \(k\) 为域，\(A\) 为结合含幺 \(k\) 代数，\(M\) 为 \(A\) 双模，\(B_\bullet(A)\) 为双模 bar 消解。
	\begin{enumerate}
		\item 从 \(M^{\mathrm r}\otimes_{A^e}B_\bullet(A)\) 的平衡关系推导一次、二次 Hochschild 边界，说明末项为何使用左 \(A\) 作用。
		\item 取 \(A=k[x,y]\)，令 \(\sigma\) 为交换 \(x,y\) 的 \(k\) 代数自同构。在底层向量空间 \(M=A\) 上规定左作用为通常乘法、右作用为 \(m\cdot a=m\sigma(a)\)。计算 \(b(1\otimes x)\) 及 \(\operatorname{HH}_0(A,M)\)。
	\end{enumerate}
\end{optionalexercise}
\begin{solution}
	1. 在标准代表 \(m\otimes_{A^e}(1\otimes a_1\otimes\cdots\otimes a_n\otimes1)\) 上，bar 微分的首项把 \(a_1\) 移到最左因子，平衡关系将它改写成系数 \(ma_1\)；末项把 \(a_n\) 移到最右因子，因 \(m\cdot(1\otimes a_n^{\mathrm{op}})=a_nm\)，系数变成 \(a_nm\)。因此
	\[
		\begin{aligned}
			b(m\otimes a)&=ma-am,\\
			b(m\otimes a\otimes b)&=ma\otimes b-m\otimes ab+bm\otimes a.
		\end{aligned}
	\]
	二次公式的末项不能改成 \(mb\)。

	2. 右作用被 \(\sigma\) 改变，所以 \(b(1\otimes x)=\sigma(x)-x=y-x\)。一般地，\(b(m\otimes a)=m\sigma(a)-am=m(\sigma(a)-a)\)，故随着 \(m\) 遍历 \(A\)，所有一次边界的线性张成正是由 \(\sigma(a)-a\)、\(a\in A\) 生成的理想。每个差在代入 \(y=x\) 后为零，因而都在 \((x-y)\) 中；取 \(a=x\) 又得到该理想的生成元。因此
	\[
		\operatorname{HH}_0(A,M)\cong k[x,y]/(x-y)\cong k[x].
	\]
\end{solution}

\begin{optionalexercise}\label{b4:ex:nonpoisson-first-deformation}
	设 \(k\) 为特征不等于 \(2\) 的域，\(A=k[x,y,z]\)。定义交替双线性运算
	\[
		\begin{aligned}
			\beta(f,g)={}&x(f_xg_y-f_yg_x)+y(f_yg_z-f_zg_y)\\
			&+z(f_zg_x-f_xg_z),
		\end{aligned}
	\]
	其中下标表示相应偏导。证明 \(f\star g=fg+\dfrac{\hbar}{2}\beta(f,g)\pmod{\hbar^2}\) 是保持单位的一阶结合形变，但它不能延拓为模 \(\hbar^3\) 的结合形变。
\end{optionalexercise}
\begin{solution}
	对每个固定的 \(f\)，映射 \(g\mapsto\beta(f,g)\) 是若干偏导子的 \(A\) 线性组合，因而是导子；另一变量同理。记 \(\phi=\dfrac12\beta\)。把两次乘积法则代入 Hochschild 微分，得到
	\[
		\begin{aligned}
			(\delta\phi)(a,b,c)
			&=a\cdot\phi(b,c)-\phi(ab,c)+\phi(a,bc)-\phi(a,b)\cdot c\\
			&=a\cdot\phi(b,c)-a\cdot\phi(b,c)-b\cdot\phi(a,c)\\
			&\qquad+\phi(a,b)\cdot c+b\cdot\phi(a,c)-\phi(a,b)\cdot c=0.
		\end{aligned}
	\]
	偏导子零化 \(1\)，故 \(\phi\) 归一化，一次结合律成立。然而
	\[
		\beta(x,y)=x,\qquad \beta(y,z)=y,\qquad \beta(z,x)=z,
	\]
	所以
	\[
		\beta(x,\beta(y,z))+\beta(y,\beta(z,x))+\beta(z,\beta(x,y))
		=x+y+z\ne0.
	\]
	一次交换子正好为 \(\phi(a,b)-\phi(b,a)=\beta(a,b)\)。若有二阶结合延拓，\cref{b4:prop:star-first-commutator-poisson}便迫使 \(\beta\) 满足 Jacobi 恒等式，与计算矛盾。这个例子说明双导子条件足以产生一次形变，却不足以使其反对称部分成为 Poisson 括号。
\end{solution}

\begin{optionalexercise}\label{b4:ex:ordered-star-polynomials}
	设 \(k\) 为特征零域，\(A=k[q,p]\)，按\cref{b4:prop:commuting-derivation-star-product}，用 \(P=\partial_q\otimes\partial_p\) 定义星积 \(\star_{\mathrm o}\)。
	\begin{enumerate}
		\item 对非负整数 \(a,b,c,d\)，计算 \((q^ap^b)\star_{\mathrm o}(q^cp^d)\)，并求 \(q^2\star_{\mathrm o}p^2\) 与 \(p^2\star_{\mathrm o}q^2\)。
		\item 证明 \(A[\![\hbar]\!]\) 在该星积下的中心恰为 \(k[\![\hbar]\!]\)。
	\end{enumerate}
\end{optionalexercise}
\begin{solution}
	1. 对第一因子求 \(q\) 导数，对第二因子求 \(p\) 导数，得到
	\[
		(q^ap^b)\star_{\mathrm o}(q^cp^d)
		=\sum_{r=0}^{\min(a,d)}
		\hbar^r\binom ar\binom dr r!\,
		q^{a+c-r}p^{b+d-r}.
	\]
	因为 \(r\) 次导数分别给出 \(a!/(a-r)!\) 与 \(d!/(d-r)!\)，再除以星积公式中的 \(r!\)，正好得到上述系数。因此
	\[
		q^2\star_{\mathrm o}p^2=q^2p^2+4\hbar qp+2\hbar^2,
		\qquad p^2\star_{\mathrm o}q^2=p^2q^2.
	\]
	其交换子的一次项为 \(4qp\)，与标准辛括号一致；二次项则不同于正文中对称形式的星积。

	2. 对任意形式级数 \(F\)，坐标函数的高阶偏导为零，直接得到
	\[
		[q,F]_{\star_{\mathrm o}}=\hbar\partial_pF,
		\qquad [p,F]_{\star_{\mathrm o}}=-\hbar\partial_qF.
	\]
	若 \(F=\sum_{n\ge0}\hbar^nf_n\) 属于中心，两式必须为零。形式级数模中乘以 \(\hbar\) 单射，故每个 \(f_n\) 都满足 \(\partial_pf_n=\partial_qf_n=0\)。特征零使非恒定单项式的某个偏导非零，所以 \(f_n\in k\)，即 \(F\in k[\![\hbar]\!]\)。反过来，来自 \(k[\![\hbar]\!]\) 的级数由标量双线性与单位条件可知都在中心。
\end{solution}

\begin{optionalexercise}\label{b4:ex:ordered-symmetric-star-equivalence}
	设 \(k\) 为特征零域，\(A=k[q,p]\)。记 \(m:A\otimes_k A\to A\) 为原乘法。在 \(A[\![\hbar]\!]\) 上分别取
	\[
		\begin{aligned}
			f\star_{\mathrm o}g&=m\exp(\hbar\partial_q\otimes\partial_p)(f\otimes g),\\
			f\star_{\mathrm s}g&=m\exp\!\left(\frac{\hbar}{2}
			(\partial_q\otimes\partial_p-\partial_p\otimes\partial_q)\right)(f\otimes g).
		\end{aligned}
	\]
	证明连续线性同构 \(T=\exp(-\dfrac{\hbar}{2}\partial_q\partial_p)\) 满足
	\[
		T(f\star_{\mathrm o}g)=T(f)\star_{\mathrm s}T(g),
	\]
	并用 \(f=q,g=p\) 检查公式中的符号。
\end{optionalexercise}
\begin{solution}
	记 \(\Delta=\partial_q\partial_p\)、\(P=\partial_q\otimes\partial_p\)、\(P^{\mathrm{op}}=\partial_p\otimes\partial_q\)，并令 \(\Delta_1=\Delta\otimes\operatorname{id}\)、\(\Delta_2=\operatorname{id}\otimes\Delta\)。两次乘积法则给出
	\[
		\Delta(fg)=\Delta(f)\cdot g+f\cdot\Delta(g)
		+\partial_q(f)\partial_p(g)+\partial_p(f)\partial_q(g),
	\]
	即 \(\Delta m=m(\Delta_1+\Delta_2+P+P^{\mathrm{op}})\)。这些算子两两交换，逐次应用这条恒等式并展开形式指数，得到
	\[
		\begin{aligned}
			Tm\exp(\hbar P)
			&=m\exp\!\left(-\frac{\hbar}{2}
			(\Delta_1+\Delta_2+P+P^{\mathrm{op}})\right)\exp(\hbar P)\\
			&=m\exp\!\left(\frac{\hbar}{2}(P-P^{\mathrm{op}})\right)
			\exp\!\left(-\frac{\hbar}{2}(\Delta_1+\Delta_2)\right)\\
			&=m\exp\!\left(\frac{\hbar}{2}(P-P^{\mathrm{op}})\right)(T\otimes T).
		\end{aligned}
	\]
	所有等式按形式系数理解，最后的 \(T\otimes T\) 在形式级数上也按每次有限和延拓。因此得到所需星积同态。\(T\) 的逆是 \(\exp(\dfrac{\hbar}{2}\Delta)\)，它保持单位且模 \(\hbar\) 为恒等，故确为允许的形变等价。对于坐标函数，\(T(q)=q\)、\(T(p)=p\)，而
	\[
		T(q\star_{\mathrm o}p)=T(qp+\hbar)
		=qp+\frac{\hbar}{2}=q\star_{\mathrm s}p.
	\]
	这也核对了指数中的负号。两种乘法的高阶公式虽不同，却由明确的形式坐标变换联系。
\end{solution}

\begin{optionalexercise}\label{b4:ex:cyclic-field-representatives}
	设 \(k\) 为域。在 \(C_n(k)=k^{\otimes_k(n+1)}\) 中记 \(e_n=1^{\otimes(n+1)}\)。
	\begin{enumerate}
		\item 计算 \(b(e_n)\) 与 \(B(e_n)\)，并在 \(\operatorname{CC}_2(k)\)、\(\operatorname{CC}_4(k)\) 中分别找出经 \(S\)、\(S^2\) 映到 \([1]\in\operatorname{HC}_0(k)\) 的循环。
		\item 对每个 \(q\ge0\)，构造 \(\operatorname{CC}_{2q}(k)\) 中的循环 \(z_q\)，使 \(z_0=e_0\)、\(S(z_q)=z_{q-1}\)。说明该构造为什么也适用于正特征。
	\end{enumerate}
\end{optionalexercise}
\begin{solution}
	1. 正次边界是交错求和，故 \(b(e_n)=e_{n-1}\) 当 \(n>0\) 为偶数，\(b(e_n)=0\) 当 \(n\) 为奇数。又有 \(t_n(e_n)=(-1)^n e_n\)。若 \(n\) 为奇数，则 \(N_n(e_n)=0\)；若 \(n\) 为偶数，则 \(N_n(e_n)=(n+1)e_n\)，且 \(1-t_{n+1}\) 在 \(e_{n+1}\) 上作用为乘以 \(2\)。所以
	\[
		B(e_n)=
		\begin{cases}
			2(n+1)e_{n+1},&n\;\text{为偶数},\\
			0,&n\;\text{为奇数}.
		\end{cases}
	\]
	在总复形中可取
	\[
		z_1=(-2e_2,e_0),\qquad z_2=(12e_4,-2e_2,e_0).
	\]
	第一条的边界为 \(-2e_1+2e_1=0\)；第二条的两个分量边界分别为 \(12e_3-2\cdot6e_3=0\)、\(-2e_1+2e_1=0\)。删去首分量便依次得到 \(z_1,z_0\)。

	2. 写 \(z_q=(c_0e_{2q},c_1e_{2q-2},\ldots,c_qe_0)\)，从 \(c_q=1\) 开始，递推规定
	\[
		c_j=-2(2q-2j-1)c_{j+1}\qquad(0\le j<q).
	\]
	这正是第 \(j\) 个边界分量为零的条件。递推只使用整数乘法，没有除法，所以将整数映到任意特征的 \(k\) 都成立。删去首分量后，剩余系数满足 \(q-1\) 的同一递推且末项仍为 \(1\)，故 \(S(z_q)=z_{q-1}\)。由\cref{b4:prop:cyclic-hh-vanishing}，\(S^q\) 在同调上是同构，\([z_q]\) 就是映到 \([1]\) 的唯一同调类。例如特征二时所有前面的系数均为零，末分量 \(e_0\) 仍给出非零类。
\end{solution}

\begin{optionalexercise}\label{b4:ex:cyclic-laurent-polynomials}
	设 \(k\) 为特征零域，\(A=k[x,x^{-1}]\)。证明
	\[
		\operatorname{HC}_0(A)=A,\qquad
		\operatorname{HC}_n(A)\cong k\quad(n\ge1).
	\]
	明确指出一次同调中一个生成元，以及 \(S:\operatorname{HC}_2(A)\to A\) 的像。
\end{optionalexercise}
\begin{solution}
	包络代数为 \(R=A^e=k[x,x^{-1},y,y^{-1}]\)，乘法映射的核为 \((x-y)\)。由于 \(R\) 是整环，乘以 \(x-y\) 单射，因此
	\[
		0\longrightarrow R\xrightarrow{\ x-y\ }R\longrightarrow A\longrightarrow0
	\]
	是投射左 \(A^e\) 模消解。与 \(A\) 张量后，\(x-y\) 作用为零，所以 \(\operatorname{HH}_0(A)=A\)、\(\operatorname{HH}_1(A)\cong A\)，而 \(\operatorname{HH}_n(A)=0\) 对 \(n\ge2\) 成立。一次同调的微分解释为 \(\Omega_{A/k}=A\,\mathrm d x\)：导子在 \(x\) 上的值决定它，且 \(\mathrm d(x^{-1})=-x^{-2}\cdot\mathrm d x\)。

	对每个整数 \(m\)，有 \(\mathrm d(x^m)=m x^{m-1}\cdot\mathrm d x\)。特征零使每个非零整数可逆，所以 \(\mathrm dA\) 包含除 \(x^{-1}\cdot\mathrm d x\) 外的所有 Laurent 单项式微分；没有一个单项式求导能产生非零的 \(x^{-1}\cdot\mathrm d x\) 项。因此
	\[
		\ker(\mathrm d:A\to\Omega_{A/k})=k,\qquad
		\Omega_{A/k}/\mathrm dA=k\cdot\left[\frac{\mathrm d x}{x}\right].
	\]
	SBI 列的低次部分于是给出
	\[
		0\longrightarrow\operatorname{HC}_2(A)\xrightarrow{S}A
		\xrightarrow{\mathrm d}\Omega_{A/k}
		\longrightarrow\operatorname{HC}_1(A)\longrightarrow0.
	\]
	因此 \(\operatorname{HC}_2(A)\cong k\)，\(S\) 的像恰为常数，\(\operatorname{HC}_1(A)\cong k\) 由 \([\mathrm d x/x]\) 生成。对 \(n\ge3\)，相邻的 \(\operatorname{HH}_n(A)\)、\(\operatorname{HH}_{n-1}(A)\) 都为零，故 \(S:\operatorname{HC}_n(A)\to\operatorname{HC}_{n-2}(A)\) 是同构，得到全部正次计算。
\end{solution}

\begin{optionalexercise}\label{b4:ex:cyclic-product-matrices}
	设 \(k\) 为域，\(r_1,\ldots,r_s\) 为正整数，\(s\ge1\)，并令 \(A=\prod_{j=1}^sM_{r_j}(k)\)。
	\begin{enumerate}
		\item 构造乘法映射 \(A\otimes_kA\to A\) 的双模截面，计算全部 \(\operatorname{HC}_n(A)\)。
		\item 求所有满足 \(\tau(ab)=\tau(ba)\) 的 \(k\) 线性泛函 \(\tau:A\to k\)。若 \(\operatorname{char}k=p>0\) 且每个 \(r_j\) 都被 \(p\) 整除，说明为什么 \([1_A]=0\in\operatorname{HC}_0(A)\)，但 \(\operatorname{HC}_0(A)\ne0\)。
	\end{enumerate}
\end{optionalexercise}
\begin{solution}
	1. 用 \(E_{uv}^{(j)}\) 表示第 \(j\) 个因子内的矩阵单位，其余分量为零。令
	\[
		e=\sum_{j=1}^s\sum_{i=1}^{r_j}E_{i1}^{(j)}\otimes E_{1i}^{(j)}.
	\]
	同一因子内使用矩阵单位乘法，不同因子的乘积为零，可得 \(\mu(e)=1_A\) 以及 \((a\otimes1)e=e(1\otimes a)\)，其中右式按外侧双模作用理解。因此 \(\rho(a)=(a\otimes1)e\) 是双模截面，\(A\) 可分。

	交换子逐因子计算，\([A,A]=\prod_j[M_{r_j}(k),M_{r_j}(k)]\)：任一分量的交换子都可用只在该分量非零的两个元素产生。各因子的普通迹因而给出
	\[
		A/[A,A]\xrightarrow{\sim}k^s,\qquad
		[(a_1,\ldots,a_s)]\longmapsto(\operatorname{tr}(a_1),\ldots,\operatorname{tr}(a_s)).
	\]
	由\cref{b4:prop:separable-cyclic-homology}，\(\operatorname{HC}_{2q}(A)\cong k^s\)、\(\operatorname{HC}_{2q+1}(A)=0\)。

	2. 条件 \(\tau(ab)=\tau(ba)\) 等价于 \(\tau\) 消去所有交换子，所以 \(\tau\) 唯一经过 \(A/[A,A]\cong k^s\)。因此所有这样的泛函恰为
	\[
		\tau(a_1,\ldots,a_s)=\sum_{j=1}^s\lambda_j\operatorname{tr}(a_j),\qquad\lambda_j\in k.
	\]
	单位类的像为 \((r_1,\ldots,r_s)\)。在题设的正特征情形，该向量为零，但每个 \([E_{11}^{(j)}]\) 映到 \(k^s\) 的相应标准基向量，故零次同调仍非零。单位元素不能代替一组交换子商的基。
\end{solution}

\begin{optionalexercise}\label{b4:ex:cyclic-inseparable-field}
	设 \(p\) 为素数，\(k=\mathbb F_p(t)\)，\(L=k(u)\) 且 \(u^p=t\)。计算相对于 \(k\) 的 \(\operatorname{HC}_1(L)\)，并与相对于 \(L\) 自身的计算比较。
\end{optionalexercise}
\begin{solution}
	\(U^p-t\) 对素元 \(t\) 满足 Eisenstein 判据，故 \([L:k]=p\)，且 \(1,u,\ldots,u^{p-1}\) 是一组 \(k\) 基。\(\Omega_{L/k}=L\,\mathrm d u\)，而对 \(a_i\in k\)，有
	\[
		\mathrm d\!\left(\sum_{i=0}^{p-1}a_i u^i\right)
		=\sum_{i=1}^{p-1}i a_i u^{i-1}\cdot\mathrm d u.
	\]
	\(1,\ldots,p-1\) 在 \(k\) 中非零，因此 \(\mathrm dL\) 的 \(k\) 基为 \(\mathrm d u,u\cdot\mathrm d u,\ldots,u^{p-2}\cdot\mathrm d u\)。由\cref{b4:cor:cyclic-first-differentials}，
	\[
		\operatorname{HC}_1(L)\cong\Omega_{L/k}/\mathrm dL
		=k\cdot[u^{p-1}\cdot\mathrm d u]\ne0,
	\]
	其中右边表示由所列同调类张成的一维 \(k\) 向量空间。若改以 \(L\) 为基域，则这是基域自身的循环同调，所有奇数次群为零。这个例子中的 \(L\) 虽然是有限维简单 \(k\) 代数，却不能套用有限维可分代数的奇数次消失结论。
\end{solution}

\begin{optionalexercise}\label{b4:ex:mixed-quasiisomorphism-cyclic}
	设 \(k\) 为域，\((C_\bullet,b,B)\)、\((C'_\bullet,b',B')\) 为非负次数混合复形，\(f_n:C_n\to C'_n\) 满足 \(f b=b'f\) 与 \(f B=B'f\)。若 \(f\) 在底层链复形同调 \(H_n(C,b)\to H_n(C',b')\) 上均为同构，证明诱导的循环复形映射也在同调上均为同构。
\end{optionalexercise}
\begin{solution}
	在循环复形的每个分量上施加 \(f\)。两条相容等式使此映射与 \(\mathrm d(x_0,x_1,\ldots)=(b x_0+B x_1,\ldots)\) 交换，因此它是链映射，并与包含 \(I\)、投影 \(S\) 相容。

	SBI 短正合列的构造只使用混合复形的三条恒等式，所以对 \(C\)、\(C'\) 同样成立。记 \(H_n=H_n(C,b)\)、\(H'_n=H_n(C',b')\)，以及 \(\operatorname{HC}_n(C)=H_n(\operatorname{CC}_\bullet(C))\)。在零次，\(\operatorname{HC}_0(C)=H_0\)、\(\operatorname{HC}_0(C')=H'_0\)，故结论成立。设所有小于 \(n\) 的循环同调映射已经是同构，比较两条长正合列中的五项
	\[
		\operatorname{HC}_{n-1}(C)\longrightarrow H_n
		\longrightarrow\operatorname{HC}_n(C)
		\longrightarrow\operatorname{HC}_{n-2}(C)\longrightarrow H_{n-1}
	\]
	与 \(C'\) 的对应五项。第一、第四项的竖箭头由归纳假设为同构，第二、第五项由题设为同构；负次循环群均为零，也不造成例外。由\cref{b4:thm:five-lemma}，中间的 \(\operatorname{HC}_n(C)\to\operatorname{HC}_n(C')\) 为同构。归纳完成证明。
\end{solution}
